\documentclass[../main.tex]{subfiles}
\begin{document}
\section{Introduction}
\label{sec:classification_intro}

Image Classification is a cornerstone of computer vision, serving as the foundational task of assigning a single, predefined label to an entire image. From distinguishing between species of animals to identifying products in a warehouse, its goal is to provide a high-level understanding of an image's primary content. The methodologies developed for image classification are often the building blocks for more complex vision tasks, such as object detection and segmentation, making its study essential for understanding the broader field.

While general-purpose classification is a well-studied area, the stakes are significantly elevated in safety-critical domains. One such domain is autonomous driving, where the correct identification of road signs is non-negotiable. A model's failure to distinguish a `Stop` sign from a `Speed Limit 60` sign could have catastrophic consequences. To ground our investigation in this high-stakes context, our experiments in this chapter will focus on the \textbf{German Traffic Sign Recognition Benchmark (GTSRB)} dataset. This dataset is a real-world, multi-class benchmark that presents a realistic challenge for modern classifiers and serves as an ideal testbed for evaluating security vulnerabilities.

The model under adversarial investigation is \textbf{EfficientNet-B1}. Unlike traditional architectures that arbitrarily scale dimensions like network depth or width, the EfficientNet family of models employs a principled method of \textit{compound scaling}, which systematically and uniformly scales the network's depth, width, and resolution. This approach has allowed EfficientNet models to achieve state-of-the-art accuracy while being significantly more computationally efficient than their predecessors. The choice of EfficientNet-B1 is deliberate: its combination of high accuracy and efficiency makes it a practical and powerful model for real-world deployment, and therefore, a crucial target for robustness analysis. Our trained EfficientNet-B1 model achieves a baseline test accuracy of \textbf{96.7\%} on the clean GTSRB test set, establishing a high-performance benchmark from which to measure the impact of adversarial attacks. If such an advanced and efficient model can be easily fooled, it underscores the severity of the adversarial threat.

In this chapter, we will therefore probe the vulnerabilities of the highly capable EfficientNet-B1 model when trained on the safety-critical GTSRB dataset. The subsequent sections will detail the specific adversarial attacks employed to generate malicious inputs and evaluate various defense strategies aimed at fortifying the model against such threats.
\section{Attacks}
\subsection{White Box Attacks}
\subsubsection{Fast Gradient Sign Method (FGSM)}

\paragraph{Introduction}
The Fast Gradient Sign Method (FGSM) is one of the earliest and most influential approaches for crafting adversarial examples against deep neural networks. Introduced by Goodfellow et al.\ in 2015, FGSM demonstrated that small, carefully chosen perturbations—imperceptible to humans—can reliably force state-of-the-art classifiers into making incorrect predictions. This work highlighted a fundamental vulnerability in high-dimensional linear behavior of models and sparked a wealth of research into both stronger attacks and robust defense mechanisms \cite{goodfellow2015explaining}.  

In their seminal paper, Goodfellow and colleagues showed that even a single-step perturbation computed from the model’s loss gradient suffices to degrade accuracy dramatically on datasets such as MNIST and CIFAR-10. The simplicity and effectiveness of FGSM make it a canonical benchmark for evaluating adversarial robustness and for developing more advanced multi-step attacks and defenses \cite{goodfellow2015explaining}.
\paragraph{Methodology}
\subparagraph{Attack Type}
The Fast Gradient Sign Method (FGSM) is a \textbf{white-box}, \textbf{untargeted} attack. As a white-box method, it requires full access to the model’s parameters and gradients. It is untargeted in the sense that its goal is simply to cause misclassification, not to steer the model toward a specific incorrect label. The adversarial example is generated to increase the loss with respect to the true class, thereby pushing the prediction away from the correct category.

\subparagraph{Parameters}
\begin{description}
  \item[$x$] The clean input image (e.g., a batch of images) to be perturbed.
  \item[$y$] The true label (or labels) associated with $x$.
  \item[$\theta$] The fixed parameters of the target neural network model.
  \item[$\epsilon$] The maximum perturbation magnitude (controls attack strength).
  \item[$J(\theta, x, y)$] The loss function used during model training (e.g., cross‐entropy).
\end{description}

\subparagraph{Algorithm Overview}
FGSM generates an adversarial example by taking a single, maximally disruptive step in the direction that increases the model’s loss most rapidly. Concretely, we compute the gradient of the loss $J(\theta, x, y)$ with respect to the input $x$, extract its sign to form a unit‐step “direction,” scale that by $\epsilon$, and add it to the original image. Finally, we clip the result to the valid pixel range to ensure a legitimate image. This one‐step procedure is efficient and often sufficient to cause high‐confidence misclassifications, making it a standard baseline for evaluating robustness.

\begin{breakablealgorithm}{Fast Gradient Sign Method (FGSM) Attack}
\begin{algorithmic}[1]
\Require $x$, $y$, $\theta$, $\epsilon$
\Ensure $x_{\mathrm{adv}}$
\State Enable gradient tracking on $x$ 
\State Compute loss: $L \gets J(\theta, x, y)$
\State Compute gradient: $g \gets \nabla_x L$ via backpropagation
\State Form perturbation: $\delta \gets \epsilon \cdot \mathrm{sign}(g)$
\State Generate adversarial example: $x_{\mathrm{adv}} \gets x + \delta$
\State Clip to valid range: $x_{\mathrm{adv}} \gets \mathrm{clip}(x_{\mathrm{adv}}, 0, 1)$
\State \Return $x_{\mathrm{adv}}$
\end{algorithmic}
\end{breakablealgorithm}

\begin{figure}[h!] % [h!] tries to place it "here" forcefully. Use [htbp] for more flexibility.
  \centering % Center the whole figure content

  % --- First Image: Original Traffic Sign ---
  \begin{subfigure}[b]{0.28\textwidth}
    \centering
    \includegraphics[width=0.7\textwidth]{images/CV/Classification/20.jpg} % YOUR ORIGINAL IMAGE
    \caption*{$x$\\\textit{"speed limit 20"}\\71\% confidence}
    % \label{fig:sub-original-traffic} % Optional label for subfigure
  \end{subfigure}
  \hfill % Distribute horizontal space
  % --- Plus and Multiplier Operator ---
  \begin{subfigure}[b]{0.05\textwidth} % Width for the operator's subfigure
    \centering
     \vspace{0.5cm} % Removed, relying on [b] alignment for all subfigures
    $+\,0.007\,\times$
  \end{subfigure}
  \hfill % Distribute horizontal space
  % --- Second Image: Noise Pattern ---
  \begin{subfigure}[b]{0.28\textwidth}
    \centering
    \includegraphics[width=0.7\textwidth]{images/CV/Classification/noise_20.png} % YOUR NOISE IMAGE
    % Caption adjusted as per "remove the class for gradient and confidence":
    \caption*{$\mathrm{sign}(\nabla_x J(\theta, x, y))$\\
    \\ \hfill
    \\ \hfill
}
    % \label{fig:sub-noise-traffic} % Optional label for subfigure
  \end{subfigure}
  \hfill % Distribute horizontal space
  % --- Equals Sign Operator ---
  \begin{subfigure}[b]{0.05\textwidth} % Width for the operator's subfigure
    \centering
    % \vspace{1.2cm} % Removed, relying on [b] alignment
    $=$
  \end{subfigure}
  \hfill % Distribute horizontal space
  % --- Third Image: Adversarial Traffic Sign ---
  \begin{subfigure}[b]{0.28\textwidth}
    \centering
    \includegraphics[width=0.7\textwidth]{images/CV/Classification/Adv_20_FGSM_007.png} % YOUR ADVERSARIAL IMAGE
    \caption*{$x + \epsilon \cdot \mathrm{sign}(\nabla_x J(\theta, x, y))$\\\textit{"Slippery road"}\\77\% confidence}
    % \label{fig:sub-adversarial-traffic} % Optional label for subfigure
  \end{subfigure}

  \caption{Visual illustration of the FGSM attack on a traffic sign. A small, imperceptible perturbation, crafted using the gradient of the loss function, causes the model to misclassify the image with high confidence.}
\label{fig:fgsm-visual-traffic-sign} % Use a descriptive and unique label
\end{figure}

Figure~\ref{fig:fgsm-visual-traffic-sign} provides a clear visual demonstration of the FGSM attack. The process begins with a clean input image, $x$, correctly classified by the model as a ``speed limit 20'' sign with 71\% confidence. As detailed in the algorithm, the attack computes the gradient of the loss function with respect to the input pixels and takes its sign to determine the direction of the perturbation. This direction, $\text{sign}(\nabla_x J(\theta, x, y))$, is then scaled by a small magnitude, $\epsilon$ (here, 0.007), to form the final perturbation. When this seemingly random, low-magnitude noise is added to the original image, the resulting adversarial example, $x_{\text{adv}}$, remains visually almost identical to the original. However, it successfully deceives the model, causing a high-confidence misclassification to ``slippery road'' with 77\% confidence, demonstrating the efficacy of the attack.

\paragraph{Results}
The model's robustness was evaluated against the FGSM attack on a random subset of 512 test images. While the model achieved a baseline accuracy of 96\% on unperturbed data, its performance degraded significantly as the adversarial perturbation magnitude ($\epsilon$) increased.

The quantitative impact of the attack is detailed in Table~\ref{tab:fgsm_results_Classification}. Even a minimal perturbation of $\epsilon=0.007$ was sufficient to cause a 51\% drop in accuracy. At the highest tested magnitude of $\epsilon=0.1$, the model's accuracy collapsed to just 9.3\%, rendering it almost entirely ineffective.

Figure~\ref{fig:fgsm_examples} provides a qualitative illustration of this vulnerability. With a small $\epsilon$ of 0.01, the perturbations are nearly imperceptible to the human eye, yet they are potent enough to cause misclassifications (e.g., changing 'No passing' to 'Vehicles over 3.5 tons prohibited'). When $\epsilon$ is increased to 0.1, the adversarial noise becomes visually apparent, and the attack consistently forces the model to make high-confidence errors on all examples shown.

\begin{figure}[H] % Using [H] to place it here, as in the original
    \centering
    \begin{subfigure}[b]{0.45\textwidth}
        \centering
        \includegraphics[width=\textwidth]{images/CV/Classification/FGSM_0_01.png} 
        % No label needed for subfigure if not referenced individually
    \end{subfigure}
    \begin{subfigure}[b]{0.45\textwidth}
        \centering
        \includegraphics[width=\textwidth]{images/CV/Classification/FGSM_0_1.png} 
        % No label needed for subfigure if not referenced individually
    \end{subfigure}
    \caption{Visual examples of FGSM attacks. Left: A small perturbation ($\epsilon=0.01$) results in subtle image changes but can still cause misclassifications. Right: A larger perturbation ($\epsilon=0.1$) introduces visible noise and consistently deceives the model.}
    \label{fig:fgsm_examples}
\end{figure}

\begin{table}[H]
    \centering
    \footnotesize
    \caption{Model accuracy under FGSM attack with varying perturbation magnitudes ($\epsilon$). Results are evaluated on 512 test samples, with a baseline accuracy of 96\% on clean data.}
    \label{tab:fgsm_results_Classification}
    \begin{tabular}{@{\hspace{0.3cm}}cS[table-format=1.3]S[table-format=2.2]S[table-format=2.2]@{\hspace{0.3cm}}}
    \toprule
    {$\epsilon$} & {Accuracy (\%)} & {Accuracy Drop (\%)} \\
    \midrule
    0.007 & 44.96 & 51.04 \\
    0.010 & 34.30 & 61.70 \\
    0.030 & 12.79 & 83.21 \\
    0.100 &  9.30 & 86.70 \\
    \bottomrule
    \end{tabular}
\end{table}

\subsubsection{Projected Gradient Descent (PGD)}

\paragraph{Introduction}

Projected Gradient Descent (PGD) is a powerful iterative attack method introduced by Madry et al.\ in their work on adversarial robustness \cite{madry2018towardspgd}. It is widely regarded as one of the most effective first-order adversarial attacks, capable of producing stronger perturbations than single-step methods like FGSM by applying multiple small gradient steps.

PGD operates by iteratively perturbing the input in the direction of the gradient of the loss function and projecting the perturbed input back into an $\epsilon$-bounded $\ell_\infty$ ball around the original input. This ensures that the adversarial example remains within a specified perturbation budget. Due to its effectiveness and reliability, PGD is often considered the “universal first-order adversary” and is commonly used both as a benchmark attack and as a training strategy (i.e., adversarial training) to improve model robustness \cite{madry2018towardspgd}.

\paragraph{Methodology}

\subparagraph{Attack Type}
The Projected Gradient Descent (PGD) attack is a \textbf{white-box}, \textbf{untargeted}, and \textbf{iterative} attack. As a white-box method, it requires full access to the model’s parameters and gradients. It is untargeted in that its objective is to cause misclassification of the input without steering the prediction toward any specific incorrect label. The attack refines the basic idea of FGSM by applying multiple small perturbation steps, with each step constrained within an $\ell_\infty$-norm ball of radius $\epsilon$ centered around the original input. After each step, the perturbation is projected back into this bounded space, ensuring the adversarial example stays within the allowed perturbation budget.

\subparagraph{Parameters}
\begin{description}
  \item[$x$] The clean input image to be perturbed.
  \item[$y$] The true label associated with $x$.
  \item[$\theta$] The parameters of the target neural network model.
  \item[$\epsilon$] The maximum allowable perturbation (controls the size of the $\ell_\infty$ ball).
  \item[$\alpha$] The step size used for each gradient update.
  \item[$n$] The number of iterations (steps) to apply the attack.
  \item[$J(\theta, x, y)$] The loss function used during training (e.g., cross-entropy).
\end{description}
\subparagraph{Algorithm Overview}
The PGD attack is an iterative extension of the Fast Gradient Sign Method (FGSM). It starts by adding a small random initialization within the $\ell_\infty$ ball of radius $\epsilon$ around the original input $x$. Then, for a fixed number of iterations $n$, the method performs gradient ascent on the loss function with respect to the input, stepping in the direction that increases the loss. After each step, the adversarial example is projected back into the valid $\ell_\infty$ ball to ensure that the perturbation does not exceed the maximum allowed distortion $\epsilon$. The adversarial example is also clipped to stay within the valid input range (e.g., pixel values between 0 and 1). This iterative process allows PGD to produce much stronger adversarial examples than single-step methods like FGSM.

\begin{breakablealgorithm}{Projected Gradient Descent (PGD) Attack}
\begin{algorithmic}[1]
\Require $x$, $y$, $\theta$, $\epsilon$, $\alpha$, $n$
\Ensure $x_{\mathrm{adv}}$
\State Initialize $x_{\mathrm{adv}} \gets x + \text{Uniform}(-\epsilon, \epsilon)$
\State Clip $x_{\mathrm{adv}} \gets \mathrm{clip}(x_{\mathrm{adv}}, 0, 1)$
\For{$i = 1$ to $n$}
    \State Enable gradient tracking on $x_{\mathrm{adv}}$
    \State Compute loss: $L \gets J(\theta, x_{\mathrm{adv}}, y)$
    \State Compute gradient: $g \gets \nabla_{x_{\mathrm{adv}}} L$
    \State Update adversarial example: $x_{\mathrm{adv}} \gets x_{\mathrm{adv}} + \alpha \cdot \mathrm{sign}(g)$
    \State Project: $x_{\mathrm{adv}} \gets \mathrm{clip}(x_{\mathrm{adv}}, x - \epsilon, x + \epsilon)$
    \State Clip to valid range: $x_{\mathrm{adv}} \gets \mathrm{clip}(x_{\mathrm{adv}}, 0, 1)$
\EndFor
\State \Return $x_{\mathrm{adv}}$
\end{algorithmic}
\end{breakablealgorithm}

\paragraph{Results}

To evaluate the effectiveness of the PGD attack, we used the same \textbf{EfficientNetB1} model trained on the \textbf{German Traffic Sign Recognition Benchmark (GTSRB)} dataset, which achieved \textbf{96\% accuracy} on the clean test set. For adversarial evaluation, we applied PGD with varying perturbation budgets $\epsilon$ and corresponding step sizes $\alpha$, using a random subset of \textbf{512 test images}.

Table~\ref{tab:pgd_results} shows how model performance degrades under the PGD attack. Compared to FGSM, PGD leads to significantly lower accuracies at similar or even slightly higher $\epsilon$ values, due to its iterative nature. At $\epsilon=0.01$, accuracy already drops to 35\%. At $\epsilon=0.1$, the attack is completely successful, driving the model’s accuracy down to 0\%.

\begin{table}[H]
\centering
\footnotesize
\caption{Accuracy of PGD attacks on EfficientNetB1 evaluated on 512 GTSRB test samples using \textbf{4 iterations} per attack. The original model accuracy without attack is \textbf{96\%}.}
\label{tab:pgd_results}
\begin{tabular}{@{\hspace{0.3cm}}lS[table-format=2.2]S[table-format=2.2]c@{\hspace{0.3cm}}}
\toprule
\textbf{PGD Parameters} & {Accuracy (\%)} & {Accuracy Drop (\%)} & {Iterations} \\
\midrule
$\epsilon=0.007$, $\alpha=0.00175$  & 36.43 & 59.57 & 4 \\
$\epsilon=0.01$, $\alpha=0.0025$  & 21.71 & 74.29 & 4 \\
$\epsilon=0.03$, $\alpha=0.0075$  & 2.71  & 93.29 & 4 \\
$\epsilon=0.10$, $\alpha=0.0250$  & 0.00  & 96.00 & 4 \\
\bottomrule
\end{tabular}
\end{table}
\clearpage
\subsubsection{Momentum Iterative Fast Gradient Sign Method(MI-FGSM)}
\paragraph{Introduction}

The Momentum Iterative Fast Gradient Sign Method (MI-FGSM) evolved to address the limitations of earlier gradient-based attack methods in crafting effective adversarial examples, particularly to improve their robustness and transferability across different model scenarios. The progression began with a single-step method that used gradient perturbations to mislead models, which, while simple and efficient, was constrained by its reliance on a linear decision boundary assumption, limiting its effectiveness for larger perturbations. This led to the development of an iterative approach that applied multiple small perturbations, enhancing white-box attack success by thoroughly exploring the loss landscape. However, this method often overfit to the target model, reducing its ability to transfer to other models, a critical drawback for black-box attacks. To overcome this, MI-FGSM was introduced, incorporating the momentum concept from optimization techniques. Momentum stabilizes update directions by accumulating gradient velocities across iterations, preventing the attack from getting trapped in local maxima and improving transferability, thus balancing attack potency across both white-box and black-box scenarios.
\cite{dong2018boosting}
\paragraph{Methodology}

\subparagraph{Attack Type}
The Momentum Iterative Fast Gradient Sign Method (MI-FGSM) is a white-box adversarial attack, requiring full access to the target model's parameters and gradients for precise perturbation crafting. It is an iterative method, applying multiple small perturbation steps to refine the adversarial example, which enhances its ability to explore the loss landscape effectively. MI-FGSM operates as an untargeted attack, aiming to maximize the loss function to induce any misclassification. During its iterations, MI-FGSM updates the input by adding perturbations proportional to the sign of an accumulated momentum term, which integrates past gradients to stabilize the direction and prevent entrapment in local maxima, thereby improving attack success and transferability across models.

\subparagraph{Parameters}
\begin{itemize}
    \item $x$ The clean input image to be perturbed.
    \item $y$ The true label associated with $x$.
    \item $f$ The parameters of the target neural network model.
    \item $\epsilon$ The maximum allowable perturbation (controls the size of the $L_\infty$ ball).
    \item $T$ The number of iterations (steps) to apply the attack.
    \item $\alpha$ The step size used for each gradient update.
    \item $\mu$ The decay factor controlling the influence of past gradients (e.g., typically set to 1.0).
    \item $J(f, x, y)$ The loss function used during training (e.g., cross-entropy).
\end{itemize}
\subparagraph{Algorithm Overview}

The Momentum Iterative Fast Gradient Sign Method (MI-FGSM) is an advanced algorithm designed to generate adversarial examples by iteratively perturbing input data to maximize the loss function of a deep neural network, thereby inducing misclassification. Building on the iterative gradient-based approach, MI-FGSM enhances attack effectiveness by incorporating a momentum term that accumulates gradient information across iterations. This momentum stabilizes the perturbation direction, preventing the algorithm from being trapped in poor local maxima and improving the transferability of adversarial examples to other models. The process involves initializing the perturbation, computing gradients at each step, updating the momentum, and adjusting the input with a controlled step size, all while adhering to a predefined perturbation bound. The following pseudocode outlines the implementation.
\begin{breakablealgorithm}{Momentum Iterative Fast Gradient Sign Method (MI-FGSM)}

\begin{algorithmic}[1]
\Require $\mathbf{x}$, $y$, $f$, $\epsilon$, $\alpha$, $T$, $\mu$
\Ensure $\mathbf{x}^*$
\State Initialize $\mathbf{g}_0 \gets 0$, $\mathbf{x}_0^* \gets \mathbf{x}$
\State Set $\alpha \gets \epsilon / T$
\For{$t = 0$ to $T-1$}
    \State Compute gradient: $\nabla_{\mathbf{x}} J(\mathbf{x}_t^*, y) \gets \nabla_{\mathbf{x}} J(f(\mathbf{x}_t^*), y)$
    \State Update momentum: $\mathbf{g}_{t+1} \gets \mu \cdot \mathbf{g}_t + \frac{\nabla_{\mathbf{x}} J(\mathbf{x}_t^*, y)}{||\nabla_{\mathbf{x}} J(\mathbf{x}_t^*, y)||_1}$
    \State Update adversarial example: $\mathbf{x}_{t+1}^* \gets \mathbf{x}_t^* + \alpha \cdot \text{sign}(\mathbf{g}_{t+1})$
\EndFor
\State \Return $\mathbf{x}^* \gets \mathbf{x}_T^*$
\end{algorithmic}
\end{breakablealgorithm}



\paragraph{Results}

The model test accuracy before the attack is 96.31\%.

\begin{table}[h]
\centering
\caption{MI-FGSM Parameters}
\begin{tabular}{cccc}
\toprule
\textbf{MI-FGSM Parameters} & \textbf{Accuracy (\%)} & \textbf{Accuracy Drop (\%)} & \textbf{Iterations} \\
\midrule
$\epsilon = 0.03$, $\mu = 0.5$ & 2.52 & 93.79 & 10 \\
$\epsilon = 0.03$, $\mu = 1.0$ & 2.71 & 93.60 & 10 \\
$\epsilon = 0.01$, $\mu = 0.5$ & 20.54 & 75.77 & 10 \\
$\epsilon = 0.01$, $\mu = 1.0$ & 24.03 & 72.28 & 10 \\
\bottomrule
\end{tabular}
\end{table}

The results reveal that test accuracy is highly sensitive to the perturbation size $\epsilon$. A larger $\epsilon = 0.03$ leads to a severe accuracy drop (2.52\%–2.71\%), indicating greater attack efficacy, with a marginal increase in accuracy (0.19\%) when $\mu$ rises from 0.5 to 1.0, suggesting a stabilizing effect of momentum. Conversely, a smaller $\epsilon = 0.01$ yields higher accuracies (20.54\%–24.03\%), with an increase of 3.49\% as $\mu$ increases from 0.5 to 1.0, highlighting improved model resilience and momentum’s role in enhancing transferability. The fixed iteration count (10) ensures consistent attack strength across configurations.
\clearpage

\subsubsection{Jacobian-based Saliency Map (JSMA) }
\paragraph{Introduction} 
The Jacobian-based Saliency Map Attack (JSMA) is a targeted adversarial attack method that exposes the vulnerability of neural networks to subtle, intentional perturbations. It operates by analyzing how changes in individual input features influence the model's output, using a construct known as the forward derivative or Jacobian. This derivative captures the sensitivity of each output class to changes in input features.

JSMA generates adversarial saliency maps that identify which features, when modified, are most likely to increase the model's confidence in a specific target class while reducing confidence in all other classes. The attack modifies only a small subset of input features, typically through a greedy search strategy guided by the saliency map, making the changes both minimal and highly effective.

Unlike traditional gradient-based attacks that perturb inputs based on the overall loss function, JSMA provides fine-grained control over feature selection by focusing on the direct relationship between inputs and outputs. In practice, it has been shown to produce highly successful targeted attacks while maintaining a high level of visual similarity to the original input, as demonstrated on image classification tasks such as handwritten digit recognition.

JSMA's strength lies in its interpretability and precision, making it a powerful tool for probing model weaknesses and for understanding how specific inputs drive model decisions. It also serves as a foundation for evaluating and improving the robustness of neural network architectures in adversarial settings 
\cite{papernot2016limitations}

\paragraph{Methodology}
\subparagraph{Attack Type}



The Jacobian-based Saliency Map Attack (JSMA) is a targeted, white-box adversarial attack designed to manipulate the output of deep neural networks (DNNs) by crafting adversarial samples that induce specific misclassifications. Unlike untargeted attacks that aim to cause any misclassification, JSMA focuses on forcing a DNN to classify an input into a predefined target class chosen by the adversary. This targeted approach makes JSMA particularly effective in scenarios where specific misclassification outcomes are desired, such as bypassing security systems or altering the behavior of autonomous systems. JSMA operates in a white-box setting, requiring full knowledge of the DNN’s architecture, including its layers, activation functions, and trained parameters. By leveraging the Jacobian matrix of the DNN, JSMA constructs adversarial saliency maps that identify the input features most likely to influence the network’s output when perturbed. These maps guide the selective modification of input features, ensuring minimal changes to the original input while achieving the desired misclassification. This precision distinguishes JSMA from other adversarial attacks, such as gradient-based methods that apply broader perturbations, making it a powerful tool for exploiting DNN vulnerabilities with subtle alterations.


\subparagraph{Parameters}

\begin{itemize}
    \item $X$ The clean input sample to be perturbed.
    \item $Y^*$ The target adversarial output class chosen by the adversary.
    \item $F$ The function learned by the target neural network during training.
    \item $\gamma$ The maximum allowable distortion (controls the fraction of features perturbed, e.g., 10\% of input features).
    \item $T$ The maximum number of iterations for the perturbation process.
    \item $\theta$ The feature variation parameter (step size for perturbing selected input features).
    \item $S(X, Y^*)$ The adversarial saliency map computed using the forward derivative $J_F(X)$.
    \item $J_F(X)$ The Jacobian matrix of the function $F$ with respect to input $X$ (forward derivative).
    \item $N$ Number of Pixels to perturb at each iteration
\end{itemize}


These parameters enable JSMA to balance effectiveness and stealth, with $\theta$ and $\gamma$ ensuring minimal perturbations and $t$ directing the attack’s outcome. Proper tuning, as demonstrated in the paper’s experiments, allows JSMA to craft adversarial samples that are both highly successful and difficult to detect.


\subparagraph{Algorithm Overview}
\begin{breakablealgorithm}{Jacobian-based Saliency Map Attack (JSMA)}
\begin{algorithmic}[1]
\Require $\mathbf{X}$, $\mathbf{Y}^*$, $\mathbf{F}$, $\gamma$, $\theta$
\Ensure $\mathbf{X}^*$
\State Initialize $\mathbf{X}^* \gets \mathbf{X}$, $\delta_{\mathbf{X}} \gets \mathbf{0}$, $\mathbf{I} \gets \{1, \ldots, |\mathbf{X}|\}$
\While{$\mathbf{F}(\mathbf{X}^*) \neq \mathbf{Y}^*$ and $||\delta_{\mathbf{X}}||_0 < \gamma \cdot |\mathbf{X}|$}
    \State Compute forward derivative: $J_{\mathbf{F}}(\mathbf{X}^*) \gets \frac{\partial \mathbf{F}(\mathbf{X}^*)}{\partial \mathbf{X}^*}$
    \State Compute saliency map: $S(\mathbf{X}, \mathbf{Y}^*) \gets \text{saliency\_map}(J_{\mathbf{F}}(\mathbf{X}^*), \mathbf{I}, \mathbf{Y}^*)$
    \State Select feature: $i_{\max} \gets \arg \max_i S(\mathbf{X}, \mathbf{Y}^*)[i]$
    \State Update feature: $\mathbf{X}^*_{i_{\max}} \gets \mathbf{X}^*_{i_{\max}} + \theta$
    \State Update perturbation: $\delta_{\mathbf{X}} \gets \mathbf{X}^* - \mathbf{X}$
\EndWhile
\State \Return $\mathbf{X}^*$
\end{algorithmic}

\end{breakablealgorithm}
\paragraph{Results}
\subparagraph{Stop Sign Sample}
The following parameters were used for the attack on the Stop Sign sample(ground truth label: Class 14):
\begin{itemize}
    \item $T$: Maximum number of iterations = 15
    \item $Y^*$: Target adversarial output class = Class 17 ``No Entry''
    \item $\theta$: Feature variation parameter = 255 (unnormalized images were used)
    \item $\gamma$: Maximum allowable distortion = 20\%
    \item $N$: Number of pixels to perturb at each iteration = 50
\end{itemize}

Before the attack, the clean Stop Sign image was classified as ``Stop'' with a confidence of 92.36\%. After the attack, the adversarial Stop Sign image, generated using JSMA with minimal perturbation, was misclassified as ``No Entry'' with a confidence of 70.18\%.

\begin{figure}[ht!]
    \centering
    \begin{subfigure}[t]{0.45\textwidth}
        \includegraphics[height=5.5cm]{images/CV/Classification/JSMA/clean_stopimage_jsma_1.png}
        \caption{Clean Stop Sign image}
        \label{fig:clean_stopimage}
    \end{subfigure}
    \hfill
    \begin{subfigure}[t]{0.45\textwidth}
        \includegraphics[height=5.5cm]{images/CV/Classification/JSMA/adv_stopimage_jsma_1.png}
        \caption{Adversarial Stop Sign image generated using JSMA}
        \label{fig:adv_stopimage}
    \end{subfigure}
    \caption{Comparison of clean and adversarial Stop Sign images generated using the Jacobian-based Saliency Map Attack (JSMA)}
    \label{fig:stopimage_comparison}
\end{figure}

\subparagraph{Speed Limit 80 Sign Sample}
The following parameters were used for the attack on the Speed Limit 80 Sign sample (ground truth label: Class 5):
\begin{itemize}
    \item $T$: Maximum number of iterations = 15
    \item $Y^*$: Target adversarial output class = Class 2 ``Speed Limit 50''
    \item $\theta$: Feature variation parameter = 255 (unnormalized images were used)
    \item $\gamma$: Maximum allowable distortion = 10\%
    \item $N$: Number of pixels to perturb at each iteration = 100
\end{itemize}

Before the attack, the clean Speed Limit 80 Sign image was classified as ``Speed Limit 80'' with a confidence of 95.54\%. After the attack, the adversarial Speed Limit 80 Sign image, generated using JSMA with minimal perturbation, was misclassified as ``Speed Limit 50'' with a confidence of 84.07\%.

\begin{figure}[htbp]
    \centering
    \begin{subfigure}[t]{0.45\textwidth}
        \includegraphics[height=5.5cm]{images/CV/Classification/JSMA/clean_speed80_jsma_2.png}
        \caption{Clean Speed Limit 80 Sign image}
        \label{fig:clean_speed80}
    \end{subfigure}
    \hfill
    \begin{subfigure}[t]{0.45\textwidth}
        \includegraphics[height=5.5cm]{images/CV/Classification/JSMA/adv_speed80_jsma_2.png}
        \caption{Adversarial Speed Limit 80 Sign image generated using JSMA}
        \label{fig:adv_speed80}
    \end{subfigure}
    \caption{Comparison of clean and adversarial Speed Limit 80 Sign images generated using the Jacobian-based Saliency Map Attack (JSMA)}
    \label{fig:speed80_comparison}
\end{figure}
\clearpage
\subparagraph{Vechiles over 3.5 tons prohibited}
The following parameters were used for the attack on the Speed Limit 80 Sign sample (ground truth label: Class 16):
\begin{itemize}
    \item $T$: Maximum number of iterations = 15
    \item $Y^*$: Target adversarial output class = Class 10 ``No passing for vehicles over 3.5 tons''
    \item $\theta$: Feature variation parameter = 255 (unnormalized images were used)
    \item $\gamma$: Maximum allowable distortion = 10\%
    \item $N$: Number of pixels to perturb at each iteration = 100
\end{itemize}

Before the attack, Clean  image classified as "Vehicles over 3.5 tons prohibited" with confidence 89.57\%. After the attack, Adversarial image generated using JSMA with minimal perturbation, misclassified as "No passing for vehicles over 3.5 tons" with confidence 87.21\%
\begin{figure}[htbp]
    \centering
    \begin{subfigure}[t]{0.45\textwidth}
        \includegraphics[height=5.5cm]{images/CV/Classification/JSMA/clean_vehicleover3.5tons_jsma_3.png}
        \caption{Clean  image classified as "Vehicles over 3.5 tons prohibited" }
        \label{fig:adv_stopimage_2}
    \end{subfigure}
    \hfill
    \begin{subfigure}[t]{0.45\textwidth}
        \includegraphics[height=5.5cm]{images/CV/Classification/JSMA/adv_vehicleover3.5tons_jsma_3.png}
        \caption{Adversarial image generated using JSMA }
        \label{fig:adv_stopimage_1}
    \end{subfigure}
    \caption{Comparison of clean and adversarial example of Class Vechiles over 3.5 Tons images generated using the Jacobian-based Saliency Map Attack (JSMA) }
    \label{fig:stopimage_comparison}
\end{figure}

\subparagraph{JSMA computed for whole Dataset}  

Since JSMA is a targeted attack, a worst-case analysis was employed to evaluate its impact on the test set. In this evaluation, the attacker attempts to misclassify each input into a single, fixed target class, regardless of the input’s true label. This setting reflects a highly constrained but severe adversarial scenario, commonly used to benchmark the effectiveness of targeted attacks.

The attack was applied to a carefully balanced subset of the test set, containing \textbf{2 samples per class}, totaling \textbf{86 samples}. The fixed target class selected for this evaluation was \textbf{Class 0 (``Speed Limit 20'')}, which was the most frequently assigned class following the attack.

\vspace{1em}
\noindent The quantitative results of the JSMA attack are summarized in Table~\ref{tab:jsma-results}.
\begin{table}[H]
\centering
\begin{tabular}{ll}
\toprule
\textbf{Metric} & \textbf{Value} \\
\midrule
Post-Attack Test Accuracy & 34.52\% \\
Attack Success Rate (to class 0) & 48.81\% \\
Average Distortion & 0.26\% of pixels \\
\bottomrule
\end{tabular}
\caption{JSMA targeted attack results (fixed target: Class 0) on a balanced test subset (2 samples per class, 86 total).}
\label{tab:jsma-results}
\end{table}



\vspace{1em}
\noindent
As shown in Table 6, The JSMA targeted attack was evaluated in a worst-case setting where all test samples were forced toward a fixed target class (Class 0). The attack achieved a \textbf{48.81\%} targeted success rate while preserving the original prediction in \textbf{34.52\%} of cases, all with minimal average distortion (\textbf{0.26\%}). These results confirm the effectiveness of JSMA as a precise and stealthy adversarial technique, even under severely constrained target conditions.


\subsubsection{Universal Adversarial Perturbations (UAP)}

\paragraph{Introduction}
While most adversarial attacks craft a unique, image-specific perturbation for each input, the concept of a Universal Adversarial Perturbation (UAP) represents a more general and potent class of threat. First discovered by Moosavi-Dezfooli et al. \cite{moosavidezfooli2017universaladversarialperturbations}, a UAP is a single, fixed, image-agnostic perturbation vector that, when added to a large number of different images from a dataset, can cause a model to misclassify them with high probability. The existence of such universal vectors is a remarkable finding, revealing systemic vulnerabilities and shared geometric correlations in the high-dimensional decision boundaries of deep neural networks. Instead of exploiting a weakness specific to one input, a UAP targets a common "blind spot" in the model's learned feature space. This property makes UAPs particularly menacing from a security standpoint; an attacker can pre-compute a single perturbation and deploy it to disrupt a wide range of inputs without needing to perform on-the-fly calculations. Furthermore, UAPs have demonstrated significant cross-model transferability, posing a severe threat even in black-box scenarios.

      

\paragraph{Methodology}

\subparagraph{Attack Type}
The generation of a Universal Adversarial Perturbation (UAP) is a \textbf{white-box}, \textbf{untargeted}, and \textbf{iterative} process. It is considered white-box because it requires access to the model's architecture and gradients over a dataset of images to compute the universal vector. The goal is untargeted: to cause misclassification to \textit{any} incorrect class. Once computed, the UAP itself can be deployed in a black-box manner, as it is a fixed vector that can be added to any new image.

\subparagraph{Parameters}
\begin{description}
  \item[$X$] A dataset of clean training images $\{x_1, x_2, ..., x_N\}$ used to generate the UAP.
  \item[$\theta$] The parameters of the target neural network model.
  \item[$v$] The universal perturbation vector, treated as a trainable variable.
  \item[$\epsilon$] The maximum perturbation magnitude for each pixel, enforcing an $\ell_\infty$ norm constraint.
  \item[\texttt{num\_epochs}] The number of training epochs to refine the perturbation vector.
  \item[$J(\theta, x, y)$] The loss function (e.g., cross-entropy).
\end{description}

\subparagraph{Algorithm Overview}
Instead of finding minimal updates one image at a time, this method formulates the UAP generation as a single optimization problem. The universal perturbation vector, $v$, is initialized as a zero tensor and is treated as the sole trainable parameter. The algorithm's objective is to find a single vector $v$ that, when added to any image from the dataset, maximizes the model's loss, thereby inducing misclassification.

To achieve this, the algorithm iterates through the dataset for a fixed number of epochs. In each training step, a batch of images is perturbed by adding the current universal vector $v$. The model's loss is then calculated on these perturbed images. Crucially, to maximize the loss using standard optimizers (which are designed for minimization), we optimize the \textit{negative} of the loss function. The gradient of this negative loss is computed with respect to the perturbation vector $v$, and an optimizer, such as Adam, updates $v$ in the direction of steepest ascent. After each update, the values in $v$ are clipped to the range $[-\epsilon, \epsilon]$, enforcing the $\ell_\infty$ norm constraint. This process iteratively refines a single perturbation that exploits common vulnerabilities across the entire dataset.

\begin{breakablealgorithm}{Universal Adversarial Perturbation (UAP) Generation}
\begin{algorithmic}[1]
\Require Data $X$, Target model $\mathcal{M}$, Perturbation budget $\epsilon$, Number of epochs
\Ensure Universal perturbation $v$
\State Initialize $v \gets 0$ as a trainable variable
\State Initialize an optimizer (e.g., Adam) for $v$
\For{$epoch = 1$ \textbf{to} \texttt{num\_epochs}}
    \For{each batch $(x_{\mathrm{batch}}, y_{\mathrm{batch}})$ in $X$}
        \State Compute loss gradient with respect to $v$:
        \State \quad $x_{\mathrm{adv}} \gets \text{clip}(x_{\mathrm{batch}} + v, 0, 1)$
        \State \quad $L \gets -J(\mathcal{M}(x_{\mathrm{adv}}), y_{\mathrm{batch}})$ \Comment{Negative loss for maximization}
        \State \quad $g \gets \nabla_v L$
        \State Update perturbation using the optimizer: $v \gets \text{OptimizerStep}(v, g)$
        \State Project $v$ into $\ell_\infty$ ball: $v \gets \text{clip}(v, -\epsilon, \epsilon)$
    \EndFor
\EndFor
\State \Return $v$
\end{algorithmic}
\end{breakablealgorithm}

\paragraph{Results}
To evaluate the threat of universal, image-agnostic perturbations, we generated a UAP against our EfficientNet-B1 model using the GTSRB test set. Following the methodology of the original UAP paper by Moosavi-Dezfooli et al. \cite{moosavidezfooli2017universaladversarialperturbations}, we conducted experiments across a range of $\ell_\infty$-norm constraints ($\epsilon$), operating entirely on the [0, 255] pixel scale.

The results, summarized in Table~\ref{tab:uap_multi_epsilon_results}, demonstrate a clear and direct relationship between the perturbation budget and the attack's success. As the maximum allowed per-pixel change ($\epsilon$) was increased from 5.0 to 15.0, the Fooling Rate rose significantly from 35.07\% to 77.15\%. This validates the core principle that a larger perturbation budget allows the optimizer to find a more potent universal weakness in the model.

At an $\epsilon$ of 10.0, a key benchmark used in the original paper, our attack achieved a Fooling Rate of 59.92\%. This result is highly plausible, confirming the vulnerability of our modern EfficientNetB1 architecture, while the slight difference from the ~75-85\% rates reported on older models like VGG and GoogLeNet may suggest an increased inherent robustness in EfficientNet. The high L2 norms further indicate that the attack consistently converged to a high-energy, dense noise pattern as the most effective strategy for this model. This experiment successfully confirms that our model is highly vulnerable to UAPs and that the attack's potency scales directly with the allowed perturbation magnitude.

\begin{table}[H]
    \centering
    \caption{UAP Performance Across Varying Epsilon Constraints (on 0-255 scale).}
    \label{tab:uap_multi_epsilon_results}
    \begin{tabular}{@{\hspace{0.3cm}}ccc@{\hspace{0.3cm}}}
    \toprule
    \textbf{Epsilon ($\ell_\infty$ Constraint)} & \textbf{Fooling Rate (\%)} & \textbf{L2 Norm} \\
    \midrule
     5.0 & 35.07 & 2020.35 \\
    10.0 & 59.92 & 3991.16 \\
    15.0 & 77.15 & 5919.11 \\
    \bottomrule
    \end{tabular}
\end{table}
\clearpage

\subsubsection{DeepFool}

\paragraph{Introduction}
The DeepFool attack, introduced by Moosavi-Dezfooli et al. in their 2016 paper \cite{moosavidezfooli2016deepfoolsimpleaccuratemethod}, presents an elegant, geometric approach to finding minimal adversarial perturbations. Unlike optimization-based methods like C\&W that solve a complex loss minimization problem, or single-step methods like FGSM that provide a coarse approximation, DeepFool operates on a simpler, iterative principle. It conceptualizes the classifier's decision boundaries as a set of separating hyperplanes. In each iteration, the algorithm linearizes the model at the current input's location and calculates the smallest step required to project the input directly onto the nearest decision boundary. This process is repeated until the image's classification changes. The result is a highly efficient algorithm that finds adversarial perturbations with significantly smaller norms than FGSM, often rivaling the quality of the much more computationally expensive C\&W attack. By providing an accurate yet efficient way to measure the minimal distance to a decision boundary, DeepFool serves as a crucial tool for reliably quantifying and comparing the adversarial robustness of different neural network architectures.

\paragraph{Methodology}

\subparagraph{Attack Type}
The DeepFool attack is a \textbf{white-box}, \textbf{untargeted}, and \textbf{iterative} attack. It is considered white-box as it requires access to the model's gradients and logits (pre-softmax outputs) to perform its calculations. The goal is untargeted: to find the minimal perturbation that pushes an image just across the nearest decision boundary, thus changing its classification to any other class. It is iterative because it refines the perturbation by taking small, calculated steps until a misclassification is achieved.

\subparagraph{Parameters}
\begin{description}
  \item[$x$] The clean input image to be perturbed.
  \item[$\theta$] The parameters of the target neural network model.
  \item[\texttt{max\_iter}] The maximum number of iterations for the attack loop.
  \item[\texttt{overshoot}] A small multiplicative factor applied to the final perturbation to ensure the adversarial example robustly crosses the decision boundary, rather than landing exactly on it.
\end{description}

\subparagraph{Algorithm Overview}
DeepFool operates on a geometric principle: it assumes the classifier partitions the input space with decision boundaries, and it seeks the shortest path to cross one of these boundaries. It is an iterative algorithm that, at each step, approximates the complex, non-linear classifier with a simple linear one around the current image location.

The algorithm begins with the original image $x_0$. In each iteration $i$, it performs the following steps:
\begin{enumerate}
    \item It computes the model's output logits and gradients for the current image, $x_i$.
    \item For every incorrect class $k$, it linearizes the decision boundary between the original class and class $k$. It then calculates the minimal perturbation, $\mathbf{r}_k$, needed to push $x_i$ onto this linearized boundary.
    \item It selects the overall minimal perturbation from all candidate classes—that is, the step that reaches the closest of all linearized boundaries.
    \item This minimal perturbation is added to the total accumulated perturbation.
\end{enumerate}
This loop continues until the image's predicted class changes or the maximum number of iterations is reached. Finally, a small `overshoot` factor is applied to the total perturbation to push the image firmly into the adversarial region. The result is a highly efficient method for finding small, L2-norm-minimal adversarial perturbations.

\begin{breakablealgorithm}{DeepFool Attack}
\begin{algorithmic}[1]
\Require Clean image $x_0$, Target model $\mathcal{M}$, Parameters \texttt{max\_iter}, \texttt{overshoot}
\Ensure Adversarial image $x_{\text{adv}}$
\State $x_i \gets x_0$
\State $y_{\text{orig}} \gets \text{argmax}(\mathcal{M}(x_0))$
\State $r_{\text{total}} \gets 0$
\For{$i = 1$ to \texttt{max\_iter}}
    \State $Z \gets \mathcal{M}(x_i)$ \Comment{Get logits}
    \If{$\text{argmax}(Z) \ne y_{\text{orig}}$}
        \State \textbf{break} \Comment{Attack successful}
    \EndIf
    \State Compute gradients of all logits w.r.t. $x_i$: $\nabla_{x_i} Z$
    \State $r_i \gets \infty$
    \For{each class $k \ne y_{\text{orig}}$}
        \State $w_k' \gets \nabla_{x_i} Z_k - \nabla_{x_i} Z_{y_{\text{orig}}}$ \Comment{Difference of gradients}
        \State $f_k' \gets Z_k - Z_{y_{\text{orig}}}$ \Comment{Difference of logits}
        \State Perturbation for class $k$: $r_k' \gets \frac{|f_k'|}{\|w_k'\|_2^2} w_k'$
        \If{$\|r_k'\|_2 < \|r_i\|_2$}
            \State $r_i \gets r_k'$ \Comment{Find the closest boundary}
        \EndIf
    \EndFor
    \State $r_{\text{total}} \gets r_{\text{total}} + r_i$
    \State $x_i \gets x_0 + r_{\text{total}}$ \Comment{Update image based on total perturbation from start}
\EndFor
\State $x_{\text{adv}} \gets \text{clip}(x_0 + (1 + \text{overshoot}) \cdot r_{\text{total}})$
\State \Return $x_{\text{adv}}$
\end{algorithmic}
\end{breakablealgorithm}
\paragraph{Results}

\begin{table}[H]
    \centering
    \caption{Model accuracy before and after the DeepFool Attack.}
    \label{tab:deep_fool_accuracy}
    \begin{tabular}{lcc}
        \toprule
        \textbf{Attack Method} & \textbf{Baseline Accuracy (\%)} & \textbf{Accuracy after Attack (\%)} \\
        \midrule
        DeepFool & 96.0 & 52.71 \\
        \bottomrule
    \end{tabular}
\end{table}

\noindent The DeepFool attack significantly reduced the model's accuracy from 96.0\% to 52.71\%, highlighting the model's vulnerability to minimal perturbations. This demonstrates the effectiveness of DeepFool in finding decision boundary-crossing adversarial examples with high success, even on a well-performing classifier.

\subsection{Black Box Attacks}
\subsubsection{Boundary Attack}

\paragraph{Introduction}
The Boundary Attack is a decision-based adversarial attack that operates in a black-box setting, requiring access only to the model's final predicted class label. First introduced by Brendel et al.\ in 2018 \cite{brendel2018decisionbasedadversarialattacksreliable}, it is particularly powerful because it avoids reliance on gradients or confidence scores, making it suitable for attacking systems where internal model information is completely inaccessible. At its core, the attack performs an iterative random walk along the model's decision boundary to find a minimal perturbation that causes a misclassification.

\paragraph{Methodology}

The Boundary Attack is a \textbf{black-box} and \textbf{decision-based} method. A crucial aspect of this attack is how the initial adversarial example is chosen, as this defines the attack's goal. There are two primary strategies:
\begin{enumerate}
    \item \textbf{Untargeted Attack:} An initial adversarial example is created, often by refining a pre-existing attack or adding noise. The goal is simply to minimize this perturbation while ensuring the model makes \textit{any} incorrect prediction.
    \item \textbf{Targeted Attack:} The attack is initialized with a clean image from a different, specific class (e.g., initializing with a "Speed limit 70" sign to attack a "Speed limit 30" sign). The goal is then to "morph" this starting image to be as visually similar to the original target image as possible, effectively forcing the model to misclassify the target as the starting class. This makes the second strategy a natural and effective way to perform a targeted attack.
\end{enumerate}

\subparagraph{Parameters}
\begin{description}
  \item[$x$] The original clean input that is correctly classified by the model.
  \item[$x_{\mathrm{adv}}$] The starting adversarial example, which must already be misclassified.
  \item[$\mathcal{M}$] The black-box model to be attacked (returns only predicted class labels).
  \item[$\text{max\_iterations}$] The maximum number of attack iterations.
  \item[$\epsilon$] The step size controlling the exploration during the random walk.
\end{description}




\begin{breakablealgorithm}{Boundary Attack}
\label{alg:boundary_attack_core}
\begin{algorithmic}[1]
\Require $x$, $x_{\mathrm{adv}}^{(0)}$, $\mathcal{M}$, $\epsilon$, $\delta$
\Ensure $x_{\mathrm{adv}}$
\State $i \gets 0$
\While{$\|x_{\mathrm{adv}}^{(i)} - x\| > \delta$ and $i < \text{max\_iterations}$}
    \State Sample a perturbation direction $r$ from a normal distribution
    \State Normalize and scale: $r \gets \epsilon \cdot \frac{r}{\|r\|}$
    \State Propose a candidate: $x_{\mathrm{cand}} \gets x_{\mathrm{adv}}^{(i)} + r$
    \State Project toward original input: $x_{\mathrm{cand}} \gets x + (x_{\mathrm{cand}} - x) \cdot \min\left(1, \frac{\|x_{\mathrm{adv}}^{(i)} - x\|}{\|x_{\mathrm{cand}} - x\|} \right)$
    \If{$\mathcal{M}(x_{\mathrm{cand}}) \ne \mathcal{M}(x)$}
        \State Accept: $x_{\mathrm{adv}}^{(i+1)} \gets x_{\mathrm{cand}}$
    \Else
        \State Reject: $x_{\mathrm{adv}}^{(i+1)} \gets x_{\mathrm{adv}}^{(i)}$
    \EndIf
    \State $i \gets i + 1$
\EndWhile
\State \Return $x_{\mathrm{adv}}^{(i)}$
\end{algorithmic}
\end{breakablealgorithm}
\paragraph{Experimental Results}
To evaluate the Boundary Attack on our trained model, we conducted experiments demonstrating both initialization strategies.

\subparagraph{Strategy 1: Refining an Existing Adversarial Example}
First, we demonstrate the attack's ability to minimize an existing perturbation. We began with a pre-crafted adversarial example—a noisy "Speed limit 20" sign that our model misclassifies as "Slippery road" with 100\% confidence. As shown in Figure~\ref{fig:boundary_attack}, the Boundary Attack uses this image as its starting point. It then iteratively refines the image, significantly reducing the visual distortion while keeping it on the adversarial side of the decision boundary. The final boundary example is much cleaner, yet it remains misclassified, with the model's confidence dropping to 49.97\%. This shows the attack's effectiveness in taking a powerful but obvious attack and transforming it into a much stealthier one.


\begin{figure}[h!]
    \centering
    \begin{subfigure}[b]{0.3\textwidth}
        \centering
        \includegraphics[width=\textwidth]{images/CV/Classification/existing_adv.png}
        \caption{Initial adversarial example. (Conf: 100\%)}
        \label{fig:initial_adv}
    \end{subfigure}
    \hfill
    \begin{subfigure}[b]{0.3\textwidth}
        \centering
        \includegraphics[width=\textwidth]{images/CV/Classification/boundary_adv_20.png}
        \caption{Final boundary example. (Conf: 49.97\%)}
        \label{fig:boundary_adv}
    \end{subfigure}
    \caption{Refinement of an adversarial example using the Boundary Attack. The attack starts with the noisy image (a) and minimizes the perturbation to produce the much cleaner, yet still misclassified, boundary image (b). The true label for both is ``Speed limit 20''.}
    \label{fig:boundary_attack}
\end{figure}

\subparagraph{Strategy 2: Morphing from a Different Class}
The second, more semantically interesting approach is to start from a legitimate image of another class. We tasked the attack with morphing a "Speed limit 70" sign into a "Speed limit 30" sign. Figure~\ref{fig:boundary_progression} documents the visual evolution of this process, showing the state of the image at every 10 iterations. The progression clearly illustrates the "morphing" effect as the features of the "70" are gradually replaced by those of the "30".

The culmination of this process is presented in Figure~\ref{fig:boundary_comparison}. The final adversarial image after 70 iterations is visually very similar to the true "Speed limit 30" target. Despite this, the model continues to misclassify it as "Speed limit 70" with 61.38\% confidence. This result is particularly compelling as it demonstrates a semantic failure of the model, creating an adversarial example that is deceptive in a human-interpretable way.
% First figure - Progression steps
\begin{figure}[H]
    \centering
    
    \begin{subfigure}[b]{0.13\textwidth}
        \centering
        \includegraphics[width=\linewidth]{images/CV/Classification/boundary_steps/boundary_attack_iter_0.png}
        \caption*{\small (Source)}
        \label{fig:boundary_iter0}
    \end{subfigure}
    \hfill
    \begin{subfigure}[b]{0.13\textwidth}
        \centering
        \includegraphics[width=\linewidth]{images/CV/Classification/boundary_steps/boundary_attack_iter_10.png}
        \caption*{\small 10 iterations}
        \label{fig:boundary_iter10}
    \end{subfigure}
    \hfill
    \begin{subfigure}[b]{0.13\textwidth}
        \centering
        \includegraphics[width=\linewidth]{images/CV/Classification/boundary_steps/boundary_attack_iter_20.png}
        \caption*{\small 20 iterations}
        \label{fig:boundary_iter20}
    \end{subfigure}
    \hfill
    \begin{subfigure}[b]{0.13\textwidth}
        \centering
        \includegraphics[width=\linewidth]{images/CV/Classification/boundary_steps/boundary_attack_iter_30.png}
        \caption*{\small 30 iterations}
        \label{fig:boundary_iter30}
    \end{subfigure}
    \hfill
    \begin{subfigure}[b]{0.13\textwidth}
        \centering
        \includegraphics[width=\linewidth]{images/CV/Classification/boundary_steps/boundary_attack_iter_40.png}
        \caption*{\small 40 iterations}
        \label{fig:boundary_iter40}
    \end{subfigure}
    \hfill
    \begin{subfigure}[b]{0.13\textwidth}
        \centering
        \includegraphics[width=\linewidth]{images/CV/Classification/boundary_steps/boundary_attack_iter_50.png}
        \caption*{\small 50 iterations}
        \label{fig:boundary_iter50}
    \end{subfigure}
    \hfill
    \begin{subfigure}[b]{0.13\textwidth}
        \centering
        \includegraphics[width=\linewidth]{images/CV/Classification/boundary_steps/boundary_attack_iter_60.png}
        \caption*{\small 60 iterations}
        \label{fig:boundary_iter60}
    \end{subfigure}

    \caption{Visual progression of the Boundary Attack showing the iterative morphing process.}
    \label{fig:boundary_progression}
\end{figure}

% Second figure - Final results
\begin{figure}[H]
    \centering
    
    \begin{subfigure}[b]{0.3\textwidth}
        \centering
        \includegraphics[width=\linewidth]{images/CV/Classification/boundary_steps/boundary_attack_iter_70.png}
        \caption{Final Adversarial \\ (Pred: Speed limit 70 \\ Conf: 61.38\%)}
        \label{fig:boundary_final}
    \end{subfigure}
    \hfill
    \begin{subfigure}[b]{0.3\textwidth}
        \centering
        \includegraphics[width=\linewidth]{images/CV/Classification/boundary_steps/4.jpg}
        \caption{Original Target \\ (pred: Speed limit 30) \\ Conf: 100\%} 
        \label{fig:boundary_target}
    \end{subfigure}

    \caption{Comparison between the final adversarial image and the original target image.}
    \label{fig:boundary_comparison}
\end{figure}
\clearpage 
\subsubsection{HopSkipJump Attack}

\paragraph{Introduction}

The HopSkipJumpAttack, introduced by Chen et al.\ in 2019 \cite{chen2020hopskipjumpattackqueryefficientdecisionbasedattack}, represents a significant advancement in the field of decision-based black-box attacks. It is specifically designed to be highly \textbf{query-efficient}, meaning it requires far fewer interactions with the target model to find an adversarial example compared to its predecessors, such as the standard Boundary Attack. This efficiency is a critical feature for evaluating real-world systems, which may be protected by query limits or other defenses that detect an unusually high volume of requests.

The core innovation of HopSkipJumpAttack is its novel method for estimating the model's gradient direction using only the binary decision outputs (i.e., "is this image class A or not?"). While the standard Boundary Attack relies on a less-directed random walk to explore the decision boundary, HopSkipJumpAttack uses its gradient estimate to guide the search much more effectively. This allows the attack to find high-quality adversarial examples with minimal, often imperceptible, perturbations in a fraction of the time and with significantly fewer model queries.
\paragraph{Methodology}

\subparagraph{Attack Type}
The HopSkipJumpAttack is a \textbf{black-box}, \textbf{decision-based} adversarial attack. This means it treats the target model as an opaque oracle, interacting with it solely by providing input images and observing the final predicted label. It does not require access to model gradients, architecture, or confidence scores (logits). The attack can be used for both \textbf{untargeted} objectives (causing any misclassification) and \textbf{targeted} objectives (causing a specific misclassification). Its primary strength lies in its query efficiency compared to other decision-based methods.

\subparagraph{Parameters}
The algorithm's execution is governed by several key parameters:
\begin{description}
  \item[$x_{orig}$] The original, clean input image that is correctly classified by the model.
  \item[$\mathcal{M}$] The black-box target model, which returns only a final class label for a given input.
  \item[$max\_iterations$] The total number of iterations for the main attack loop.
  \item[$max\_queries$] A hard limit on the total number of times the model can be queried to ensure efficiency.
  \item[$\delta$] The perturbation size used for the Monte Carlo gradient estimation.
  \item[$B$] The batch size, or number of random directions sampled, for estimating the gradient.
\end{description}

% --- PHASE 1 ---
\subparagraph{Phase 1: Adversarial Initialization}
The first step of any decision-based attack is to find a starting point that is already in the adversarial region. Since the attack starts with only the original clean image, this is typically achieved by adding progressively larger amounts of random noise to the original image until the target model misclassifies it. Algorithm~\ref{alg:hsja_init} outlines this process.

\begin{breakablealgorithm}{Adversarial Initialization}
\label{alg:hsja_init}
\begin{algorithmic}[1]
\Require Original image $x_{orig}$, Target model $\mathcal{M}$
\Ensure An initial adversarial example $x_{adv}$
\State $y_{orig} \gets \mathcal{M}(x_{orig})$
\State $\alpha \gets \text{initial\_noise\_factor}$
\For{$i \gets 1$ \textbf{to} max\_attempts}
    \State $\text{noise} \gets \text{Uniform}(-\alpha, \alpha)$
    \State $x_{cand} \gets \text{clip}(x_{orig} + \text{noise})$
    \If{$\mathcal{M}(x_{cand}) \neq y_{orig}$}
        \State \Return $x_{cand}$
    \EndIf
    \State $\alpha \gets \alpha \times \text{growth\_factor}$
\EndFor
\State \Return \textbf{null}
\end{algorithmic}
\end{breakablealgorithm}

% --- PHASE 2 ---
\subparagraph{Phase 2: Iterative Refinement}
Once an initial adversarial point is found, the main iterative refinement loop begins. The overall structure of this loop is shown in Algorithm~\ref{alg:hsja_main}. It orchestrates the three core components of the attack—the "Hop," "Skip," and "Jump"—to progressively move the adversarial example closer to the original image. Each of these components is detailed in the subsequent algorithms.

\begin{breakablealgorithm}{HopSkipJumpAttack (Main Loop)}
\label{alg:hsja_main}
\begin{algorithmic}[1]
\Require Original image $x_{orig}$, Target model $\mathcal{M}$
\Ensure A final, refined adversarial example $x_{final}$
\State $x_{adv} \gets \text{Adversarial-Initialization}(x_{orig}, \mathcal{M})$ \Comment{See Algorithm~\ref{alg:hsja_init}}
\If{$x_{adv}$ is \textbf{null}} \Return $x_{orig}$ \EndIf
\State $x_{best} \gets x_{adv}$
\For{$t \gets 1$ \textbf{to} max\_iterations}
    \If{query budget exceeded} \textbf{break} \EndIf
    \State $\text{grad\_est} \gets \text{Estimate-Gradient-Hop}(x_{adv}, \mathcal{M}, ...)$ \Comment{See Algorithm~\ref{alg:hsja_hop}}
    \State $x_{cand} \gets \text{Geometric-Step-Search-Skip}(x_{adv}, \text{grad\_est}, ...)$ \Comment{See Algorithm~\ref{alg:hsja_skip}}
    \State $x_{adv} \gets \text{Binary-Search-Projection-Jump}(x_{cand}, x_{orig}, \mathcal{M})$ \Comment{See Algorithm~\ref{alg:hsja_jump}}
    \If{$\|x_{adv} - x_{orig}\| < \|x_{best} - x_{orig}\|$}
        \State $x_{best} \gets x_{adv}$
    \EndIf
\EndFor
\State \Return $x_{best}$
\end{algorithmic}
\end{breakablealgorithm}

% --- COMPONENT DETAILS ---
\subparagraph{Component 1: The "Hop" (Gradient Estimation)}
The first step within the main loop is to estimate the gradient direction at the current boundary point. As shown in Algorithm~\ref{alg:hsja_hop}, this is done via a Monte Carlo method. Many random vectors are sampled, and for each, the model is queried to see if adding that small perturbation causes the image to cross the decision boundary. The average of these directional outcomes provides a robust estimate of the true gradient.

\begin{breakablealgorithm}{Gradient Estimation ("Hop")}
\label{alg:hsja_hop}
\begin{algorithmic}[1]
\Require Current adversarial point $x_{adv}$, Model $\mathcal{M}$, Perturbation size $\delta$, Batch size $B$
\Ensure Estimated gradient direction $\text{grad\_est}$
\State $y_{orig} \gets \mathcal{M}(x_{orig})$
\State $\text{grad\_est} \gets \text{zero\_vector}$
\For{$b \gets 1$ \textbf{to} $B$}
    \State Sample a random direction vector $u_b$ and normalize it: $u_b \gets u_b / \|u_b\|$
    \State $x_{cand} \gets x_{adv} + \delta \cdot u_b$
    \If{$\mathcal{M}(x_{cand}) \neq y_{orig}$}
        \State $\text{grad\_est} \gets \text{grad\_est} + u_b$
    \Else
        \State $\text{grad\_est} \gets \text{grad\_est} - u_b$
    \EndIf
\EndFor
\State \Return $\text{grad\_est} / \|\text{grad\_est}\|$
\end{algorithmic}
\end{breakablealgorithm}

\subparagraph{Component 2: The "Skip" (Geometric Step Search)}
After estimating the gradient, the algorithm must take a step away from the boundary to allow room for the next "Jump." As detailed in Algorithm~\ref{alg:hsja_skip}, it performs a search with geometrically increasing step sizes along the gradient direction until it finds a point that is safely inside the adversarial region.

\begin{breakablealgorithm}{Geometric Step Search ("Skip")}
\label{alg:hsja_skip}
\begin{algorithmic}[1]
\Require Current point $x_{adv}$, Gradient estimate $\text{grad\_est}$, Model $\mathcal{M}$
\Ensure A new adversarial point $x'_{adv}$ further from the boundary
\State $y_{orig} \gets \mathcal{M}(x_{orig})$
\State $\epsilon \gets \text{initial\_step\_size}$
\Loop
    \State $x'_{adv} \gets x_{adv} + \epsilon \cdot \text{grad\_est}$
    \If{$\mathcal{M}(x'_{adv}) \neq y_{orig}$}
        \State \Return $x'_{adv}$ \Comment{Found a valid adversarial point}
    \EndIf
    \State $\epsilon \gets \epsilon \times 2$ \Comment{Increase step size and try again}
\EndLoop
\end{algorithmic}
\end{breakablealgorithm}

\subparagraph{Component 3: The "Jump" (Binary Search Projection)}
Finally, the algorithm projects the new point back onto the decision boundary. Algorithm~\ref{alg:hsja_jump} shows how this is done with a highly efficient binary search. It interpolates between the new adversarial point and the original clean image, rapidly converging on a point that lies almost exactly on the boundary, ready for the next iteration.

\begin{breakablealgorithm}{Binary Search Projection ("Jump")}
\label{alg:hsja_jump}
\begin{algorithmic}[1]
\Require Adversarial point $x_{adv}$, Original point $x_{orig}$, Model $\mathcal{M}$
\Ensure A new adversarial point $x''_{adv}$ on the boundary
\State $y_{orig} \gets \mathcal{M}(x_{orig})$
\State $low \gets 0.0$, $high \gets 1.0$, $x''_{adv} \gets x_{adv}$
\For{$i \gets 1$ \textbf{to} max\_trials}
    \State $mid \gets (low + high) / 2$
    \State $x_{cand} \gets (1 - mid) \cdot x_{orig} + mid \cdot x_{adv}$
    \If{$\mathcal{M}(x_{cand}) \neq y_{orig}$}
        \State $high \gets mid$; \quad $x''_{adv} \gets x_{cand}$
    \Else
        \State $low \gets mid$
    \EndIf
    \If{$high - low < \text{threshold}$} \textbf{break} \EndIf
\EndFor
\State \Return $x''_{adv}$
\end{algorithmic}
\end{breakablealgorithm}

    In concert, these algorithms form a complete framework for executing the HopSkipJumpAttack. The initialization phase (Algorithm~\ref{alg:hsja_init}) secures a valid starting point, after which the main loop (Algorithm~\ref{alg:hsja_main}) uses the 'Hop', 'Skip', and 'Jump' sub-routines to iteratively refine the adversarial example. The primary advantage of this multi-stage methodology is its ability to find minimally perturbed solutions with high query efficiency. Having detailed the algorithmic underpinnings of the attack, we now proceed to evaluate its practical performance against our trained model.

\paragraph{Results}
To evaluate the practical effectiveness of the HopSkipJumpAttack, we conducted an experiment on a random subset of 100 images from the GTSRB test set against our trained EfficientNetB1 model. The attack was configured with a maximum of 20 refinement iterations and a budget of 5,000 model queries per image. The entire evaluation completed in 325.09 minutes, averaging approximately 3.25 minutes per image. The attack was highly successful, achieving a 97\% Attack Success Rate (ASR) by fooling 97 of the 100 initially correct images.

To contextualize these findings, we compare them against the benchmark results reported for ImageNet in the original paper by Chen et al. \cite{chen2020hopskipjumpattackqueryefficientdecisionbasedattack}, as its image resolution is most similar to ours. Table~\ref{tab:hsja_comparison_median} presents a side-by-side view using our new median-based metrics for a more direct and robust comparison.

\begin{table}[H]
    \centering
    \caption{Comparison of HopSkipJump Attack Performance}
    \label{tab:hsja_comparison_median}
    \begin{tabular}{@{\hspace{0.3cm}}lcc@{\hspace{0.3cm}}}
    \toprule
    \textbf{Metric} & \textbf{Original Paper (ImageNet)} & \textbf{Our Results (GTSRB)} \\
    \midrule
    Model & Inception-v3 & EfficientNetB1 \\
    Query Budget & 10,000 - 20,000 & \textbf{5,000} \\
    Attack Success Rate (ASR) & $\approx$94\% & \textbf{97\%} \\
    Median L2 Norm (0-1) & $\approx$2-3 & \textbf{6.1014} \\
    Median SSIM & Not Reported & \textbf{0.8755} \\
    \bottomrule
    \end{tabular}
\end{table}

Our observed 97\% ASR is consistent with the high efficacy reported in the original study. Furthermore, the median L0 norm of 172,800 indicates the attack generates a dense perturbation, modifying nearly every pixel, which is characteristic of this method. The high median SSIM of 0.8755 strongly confirms that, for a typical successful attack, the resulting adversarial image maintains excellent perceptual similarity to the original.

The median L2 norm of 6.10, while a more accurate comparison point than our previous average, remains higher than the ~2-3 reported in the paper. This difference is primarily attributable to the experimental **query budget**. The HopSkipJump algorithm is iterative; our budget of 5,000 queries is significantly lower than the 10,000-20,000 queries often used in the original benchmarks. A lower budget halts the optimization process earlier. While this is sufficient to find a successful misclassification, it is often insufficient to refine the perturbation to its minimal possible L2 norm, resulting in a successful but suboptimal example with a larger magnitude. The different model architectures (EfficientNetB1 vs. Inception-v3) and dataset complexities (GTSRB vs. ImageNet) are also likely contributing factors to this difference in final perturbation size.

In summary, our results validate the effectiveness of HopSkipJump on a modern architecture in a safety-critical domain. The performance metrics are consistent with the literature once key experimental parameters, most notably the query budget, are taken into account.
\begin{figure}[H]
    \centering
    \includegraphics[width=0.9\linewidth]{images/CV/Classification/hopskipjump_result.png}
    \caption{Visual result of the HopSkipJump attack. The original image (left) is correctly classified as a 'Speed limit 20' sign. After the attack, the adversarial image (right) is misclassified as a 'Slippery road' sign, despite being visually similar.}
    \label{fig:hopskipjump_visual_result}
\end{figure}
\subsubsection{One-Pixel Attack}

\paragraph{Introduction}
The One-Pixel Attack, introduced by Su et al.~\cite{Su_2019}, represents an extreme and compelling form of black-box adversarial attack. It challenges the conventional understanding of adversarial perturbations by demonstrating that a highly effective attack can be launched by modifying just a single pixel in an input image. Unlike methods that distribute a small amount of noise across many pixels, this attack concentrates the entire perturbation into one location. Its success highlights a critical vulnerability in modern deep neural networks, showing that their decision boundaries can be surprisingly fragile. The attack relies on an evolutionary search algorithm, specifically Differential Evolution (DE), to efficiently search the vast space of possible single-pixel modifications to find one that causes a misclassification.

\paragraph{Methodology}

\subparagraph{Attack Type} 

The One-Pixel Attack is a \textbf{black-box} attack, requiring no information about the model's architecture, parameters, or gradients. It only needs access to the model's output probabilities (logits or confidence scores), making it applicable to a wide range of non-differentiable or otherwise protected systems. It can be executed in both \textbf{untargeted} (cause any misclassification) and \textbf{targeted} (cause a specific misclassification) modes.

\subparagraph{Parameters}
\begin{description}
    \item[$x$] The clean input image to be perturbed.
    \item[$y_{\mathrm{true}}$] The true label of the clean image (used for untargeted attacks).
    \item[$M$] The black-box target model, which returns class probabilities for a given input.
    \item[\texttt{target\_class} (optional)] The desired incorrect label for a targeted attack. If not provided, the attack runs in untargeted mode.
    \item[\texttt{population\_size}] The number of candidate modifications (the population) in each generation of the search.
    \item[\texttt{max\_generations}] The maximum number of iterations for the evolutionary algorithm to run.
\end{description}

\subparagraph{Algorithm Overview}
The One-Pixel Attack formulates the search for an adversarial pixel as an optimization problem. To solve this, it employs a powerful black-box optimization method called \textbf{Differential Evolution (DE)}~\cite{Storn1997}. DE is an evolutionary algorithm that is particularly effective at finding a global optimum in a complex search space without needing gradient information.

At its core, DE works by maintaining a "population" of candidate solutions—in this case, different single-pixel modifications $(x, y, R, G, B)$. In each generation, it creates new "trial" solutions by combining existing ones. The key idea is its mutation strategy: a new solution is generated by taking one candidate and adding a scaled version of the \textit{difference} between two other candidates. This allows the search to be guided by the diversity of the current population. Each new trial solution is then compared to its "parent," and the one that is "fitter" (i.e., brings the model's output closer to the desired outcome) survives into the next generation.

The fitness function is defined based on the attack goal. For a targeted attack, fitness is the probability of the `target\_class`. For an untargeted attack, fitness is defined as the negative probability of the true class, encouraging the model to predict anything else. This evolutionary process continues until a successful attack is found or the generation limit is reached, as detailed in Algorithm~\ref{alg:one_pixel_attack_general}.

\begin{breakablealgorithm}{One-Pixel Attack}
\label{alg:one_pixel_attack_general}
\begin{algorithmic}[1]
\Require Clean image $x$, True label $y_{\mathrm{true}}$, Target model $M$, \texttt{population\_size}, \texttt{max\_generations}, \texttt{target\_class} (optional)

\Function{ComputeFitness}{$c$, $x$, $M$, $y_{\mathrm{true}}$, \texttt{target\_class}}
    \State $x_{\mathrm{cand}} \gets \text{ApplyPixel}(x, c)$
    \State $P_{\mathrm{cand}} \gets M(x_{\mathrm{cand}})$ \Comment{Get prediction probabilities}
    \If{\texttt{target\_class} is set}
        \State \Return $P_{\mathrm{cand}}[\texttt{target\_class}]$ \Comment{Targeted: Maximize target probability}
    \Else
        \State \Return $1.0 - P_{\mathrm{cand}}[y_{\mathrm{true}}]$ \Comment{Untargeted: Minimize true class probability}
    \EndIf
\EndFunction

\State Initialize Population $P$ with \texttt{population\_size} random modifications $c = (x, y, r, g, b)$.
\For{$g = 1$ \textbf{to} \texttt{max\_generations}}
    \For{each candidate $c_{\mathrm{parent}}$ in $P$}
        \State Generate $c_{\mathrm{child}}$ using DE mutation/crossover on $P$.
        \State $\mathrm{fitness}_{\mathrm{child}} \gets \text{ComputeFitness}(c_{\mathrm{child}}, x, M, y_{\mathrm{true}}, \texttt{target\_class})$
        \State $\mathrm{fitness}_{\mathrm{parent}} \gets \text{ComputeFitness}(c_{\mathrm{parent}}, x, M, y_{\mathrm{true}}, \texttt{target\_class})$
        \If{$\mathrm{fitness}_{\mathrm{child}} > \mathrm{fitness}_{\mathrm{parent}}$}
            \State Replace $c_{\mathrm{parent}}$ with $c_{\mathrm{child}}$ in $P$.
        \EndIf
    \EndFor
    
    \State Find $c_{\mathrm{best}}$, the candidate with highest fitness in the current population $P$.
    \State $x_{\mathrm{best}} \gets \text{ApplyPixel}(x, c_{\mathrm{best}})$
    \State $y_{\mathrm{pred}} \gets \text{argmax}(M(x_{\mathrm{best}}))$ \Comment{Get predicted class}
    
    \If{(\texttt{target\_class} is set \textbf{and} $y_{\mathrm{pred}} == \texttt{target\_class}$) \textbf{or} (\texttt{target\_class} is not set \textbf{and} $y_{\mathrm{pred}} \ne y_{\mathrm{true}}$)}
        \State \Return $x_{\mathrm{best}}$ \Comment{Attack successful}
    \EndIf
\EndFor
\State \Return $x_{\mathrm{best}}$ \Comment{Return the best image found, even if unsuccessful}
\end{algorithmic}
\end{breakablealgorithm}
\paragraph{Results}
To evaluate the practical impact of the One-Pixel Attack, we applied it to our EfficientNet-B1 model. The attack was configured with a population size of 25 and a maximum of 3 generations per image, prioritizing a swift yet effective search for vulnerabilities. The results demonstrate that even this extremely constrained and inconspicuous attack can successfully deceive a high-performance classifier.

A compelling case study is presented in Figure~\ref{fig:one_pixel_attack_example}. The original image, shown in Figure~\ref{fig:one_pixel_original_30}, was correctly identified by the model as a 'Speed limit 30' sign with 99.99\% confidence. After the attack, a single pixel was altered—a change imperceptible to a human observer, This minute modification was sufficient to cause a catastrophic misclassification. The model's prediction was flipped to 'Speed limit 80' with an equally high confidence of 99.17\%. This example starkly illustrates the model's brittleness, where a minimal, localized perturbation leads to a high-confidence, critical error.

\begin{figure}[H]
    \centering
    \begin{subfigure}[t]{0.48\textwidth} % Changed [b] to [t] for top alignment
        \centering
        \includegraphics[width=0.8\textwidth]{images/CV/Classification/30.jpeg} % Replace with your image path
        \caption{Original image, correctly classified as 'Speed limit 30' with 99.99\% confidence.}
        \label{fig:one_pixel_original_30}
    \end{subfigure}
    \hfill
    \begin{subfigure}[t]{0.48\textwidth} % Changed [b] to [t] for top alignment
        \centering
        \includegraphics[width=0.8\textwidth]{images/CV/Classification/30_one_pixel.jpeg} % Replace with your image path
        \caption{Adversarial image with one modified pixel (highlighted), misclassified as 'Speed limit 80' with 99.17\% confidence.}
        \label{fig:one_pixel_attacked_30}
    \end{subfigure}
    \caption{A successful One-Pixel Attack. A single-pixel change, invisible under normal viewing, causes the model to misclassify a '30' sign as an '80' sign with high confidence.}
    \label{fig:one_pixel_attack_example}
\end{figure}

The vulnerability is not an isolated phenomenon but a systemic weakness. When evaluated across the test dataset, the attack systematically degraded the model's overall performance. As summarized in Table~\ref{tab:one_pixel_accuracy}, the model's baseline accuracy of 96\% on clean data dropped to 79\%. This 17-point reduction, achieved with an attack that modifies only 0.001\% of the pixels in a 32x32 image, underscores the profound security challenge posed by such sparse adversarial perturbations. The results confirm that even state-of-the-art models can possess decision boundaries that are highly sensitive to minuscule, targeted changes.

\begin{table}[H]
    \centering
    \caption{Model accuracy before and after the One-Pixel Attack.}
    \label{tab:one_pixel_accuracy}
    \begin{tabular}{lcc}
        \toprule
        \textbf{Attack Method} & \textbf{Baseline Accuracy (\%)} & \textbf{Accuracy after Attack (\%)} \\
        \midrule
        One-Pixel & 96.0 & 79.0 \\
        \bottomrule
    \end{tabular}
\end{table}

\subsubsection{Simba Attack}
\paragraph{Introduction}


The Simple Black-box Adversarial Attack (SimBA), introduced by Guo et al. (2019), is a pioneering approach to generating adversarial examples in the black-box setting, where attackers can only query the target model's outputs without insight into its internal structure or parameters. Renowned for its exceptional query efficiency and conceptual elegance, SimBA addresses the pressing challenge of crafting effective adversarial perturbations under strict query constraints, making it highly relevant for real-world applications, such as probing machine learning APIs like Google Cloud Vision.\cite{guo2019simple} By exploiting continuous-valued confidence scores, SimBA iteratively modifies an input image by sampling random directions from a predefined orthonormal basis, steering it toward misclassification with remarkable efficiency. The method’s adaptability is a key strength, as it operates in two distinct orthonormal bases: pixel space, which directly alters individual pixel values, and Discrete Cosine Transform (DCT) coefficients, which capitalize on low-frequency components to exploit model sensitivities more effectively. This dual-basis flexibility, paired with its robust performance in practical scenarios, establishes SimBA as a formidable tool for uncovering vulnerabilities in machine learning systems within high-dimensional spaces, setting a rigorous benchmark for future research into black-box adversarial attacks.
\cite{kumar2020black}
\paragraph{Methodology}
\subparagraph{Attack Type}

The Simple Black-box Adversarial Attack (SimBA) implemented in its untargeted variant is a query-based black-box attack, designed to operate without access to the target model's internal parameters or gradients. In this approach, the adversary relies solely on querying the model's output probabilities to craft imperceptible perturbations, aiming to mislead the model into predicting any incorrect label. SimBA iteratively adjusts the input image along random directions from an orthonormal basis, using the confidence scores obtained from these queries to guide the perturbation process, making it highly efficient within the constraints of a limited query budget.
\subparagraph{Parameters}
\begin{itemize}
    \item $x$ The clean input image to be perturbed.
    \item $y$ The true label associated with $x$.
    \item $f$ The parameters of the target neural network model, accessible only via black-box queries.
    \item $\epsilon$ The step size for each perturbation update along the orthonormal basis vector.
    \item $T$ The maximum number of iterations (steps) to apply the attack, bounding the query budget.
    \item $Q$ The set of orthonormal basis vectors (standard basis for SimBA, DCT basis for SimBA-DCT).
    \item $r$ The fraction of low-frequency DCT coefficients used in SimBA-DCT (e.g., initially set to $1/8$).
    \item $p_h(y \mid x)$ The confidence score (probability) of label $y$ for input $x$, used as the loss function.
\end{itemize}


\subparagraph{Algorithm Overview}
\begin{breakablealgorithm}{Simba Attack Algorithm}
\begin{algorithmic}[1]
\Require $\mathbf{x}$, $y$, $Q$, $\epsilon$, $T$
\Ensure $\delta$
\State Initialize $\delta \gets \mathbf{0}$, $\mathbf{p} \gets p_h(y \mid \mathbf{x})$, $\mathbf{Q}_{\text{avail}} \gets Q$
\While{$\mathbf{p}_y = \max_{y'} p_h(y' \mid \mathbf{x} + \delta)$ and $t < T$}
    \State Pick random vector without replacement: $\mathbf{q} \gets \text{random}(\mathbf{Q}_{\text{avail}})$
    \State Remove $\mathbf{q}$ from $\mathbf{Q}_{\text{avail}}$
    \For{$\alpha \in \{\epsilon, -\epsilon\}$}
        \State Compute new probability: $\mathbf{p}' \gets p_h(y \mid \mathbf{x} + \delta + \alpha \mathbf{q})$
        \If{$\mathbf{p}'_y < \mathbf{p}_y$}
            \State Update perturbation: $\delta \gets \delta + \alpha \mathbf{q}$
            \State Update probability: $\mathbf{p} \gets \mathbf{p}'$
            \State \textbf{break}
        \EndIf
    \EndFor
    \State Increment iteration: $t \gets t + 1$
\EndWhile
\State \Return $\mathbf{X}^*$
\end{algorithmic}

\end{breakablealgorithm}
\subparagraph{Results}
% Paragraph and image (non-floating)
\noindent
We conducted a comparative evaluation of the original SimBA algorithm in pixel space and its frequency-domain variant, SimBA-DCT, using the GTSRB dataset under untargeted attack settings. To ensure a fair comparison, we fixed both the step size and the maximum number of allowed queries across all experiments. The primary evaluation metric was convergence speed, defined as the number of queries required to induce a misclassification. Our results demonstrate that SimBA-DCT—leveraging the Discrete Cosine Transform (DCT) basis—achieves superior query efficiency relative to the pixel-space SimBA. This improvement is attributed to SimBA-DCT’s targeted perturbations within the low-frequency components of the input, which are shown to be highly effective despite involving only a small subset of the DCT frequency spectrum, as illustrated in the following frequency mask:

\vspace{0.5em}
\begin{center}
    \includegraphics[height=7cm, width=14cm, keepaspectratio]{images/CV/Classification/Low_freq_mask.png}
    \captionsetup{type=figure}
    \captionof{figure}{Low-Frequency Mask.}
    \label{fig:dct_frequency_mask}
\end{center}

\subparagraph{SimBA on Single Samples}
To illustrate the practical difference in query efficiency between SimBA-Pixel and SimBA-DCT, we present qualitative examples of both attack variants applied to the same input samples from the GTSRB dataset. The attacks were performed under identical hyperparameter settings to ensure a fair comparison:
\begin{itemize}\itemsep0.5em
    \item $\epsilon = 0.2$
    \item $T = 2000$
    \item $Q$: Set of orthonormal basis vectors (standard basis for SimBA-Pixel, DCT basis for SimBA-DCT)
    \item $r = \frac{1}{8}$
    \item Maximum number of queries = 5000
\end{itemize}
The goal in each case is to flip the model's prediction to a different (incorrect) class using the fewest possible model queries. The resulting adversarial examples are shown below.

\begin{figure}[ht!]
    \centering
    \begin{subfigure}[t]{0.45\textwidth}
        \includegraphics[height=5.5cm]{images/CV/Classification/stop_simba_dct.png}
        \caption{SimBA-DCT: 1550 iterations, 2320 queries}
        \label{fig:simba_dct_stop}
    \end{subfigure}
    \hfill
    \begin{subfigure}[t]{0.45\textwidth}
        \includegraphics[height=5.5cm]{images/CV/Classification/stop_simba_pixel.png}
        \caption{SimBA-Pixel:1778 iterations, 2656 queries}
        \label{fig:simba_pixel_stop}
    \end{subfigure}
    \caption{SimBA attacks on a "Stop" sign image.}
    \label{fig:simba_stop_comparison}
\end{figure}

As shown in Figure~\ref{fig:simba_stop_comparison}, both adversarial examples were originally predicted as Class 14 (Stop). After the attack, the predictions flipped to Class 17 (No Entry). Notably, SimBA-DCT achieved the misclassification with fewer queries, demonstrating faster convergence.

\begin{figure}[ht!]
    \centering
    \begin{subfigure}[t]{0.45\textwidth}
        \includegraphics[height=5.5cm]{images/CV/Classification/speedlimit120_simba_dct.png}
        \caption{SimBA-DCT: 175 iterations, 264 queries}
        \label{fig:simba_dct_speed}
    \end{subfigure}
    \hfill
    \begin{subfigure}[t]{0.45\textwidth}
        \includegraphics[height=5.5cm]{images/CV/Classification/speedlimit120_simba_pixel.png}
        \caption{SimBA-Pixel: 286 iterations, 440 queries}
        \label{fig:simba_pixel_speed}
    \end{subfigure}
    \caption{SimBA attacks on a "Speed Limit 120" sign image.}
    \label{fig:simba_speed_comparison}
\end{figure}

Similarly, in Figure~\ref{fig:simba_speed_comparison}, both variants successfully changed the model's prediction from Class 8 (Speed Limit 120) to Class 7 (Speed Limit 100). Again, the DCT-based version required significantly fewer queries, further emphasizing its improved query efficiency in practice.

% Creating the subparagraph for SimBA results
\subparagraph{SimBA on Test Set}

Due to the long running time and the high number of iterations required for a single sample, we evaluated the test accuracy using 86 samples from the test set. The results are presented in Table~\ref{tab:simba-accuracy-comparison}.

% Creating the table
\begin{table}[ht]
\centering
\caption{Test-set classification accuracies for the baseline and SimBA variants}
\label{tab:simba-accuracy-comparison}
\begin{tabular}{l c}
\toprule
\textbf{Model} & \textbf{Test Accuracy (\%)} \\
\midrule
Baseline Model         & 96.00 \\
SimBA (pixel space)    & 62.79 \\
SimBA-DCT              & 65.12 \\
\bottomrule
\end{tabular}
\end{table}

% Adding a comment on the results
The results in Table~\ref{tab:simba-accuracy-comparison} demonstrate the effectiveness of both SimBA variants in reducing the test accuracy of the baseline model, which achieves a high accuracy of 96.00\%. SimBA in pixel space reduces the test accuracy to 62.79\%, indicating a significant degradation of 33.21 percentage points. SimBA-DCT, while also effective, results in a slightly higher test accuracy of 65.12\%, corresponding to a reduction of 30.88 percentage points. The lower accuracy for SimBA in pixel space aligns with its higher success rate (98.6\% for untargeted attacks) reported in the original paper, suggesting it successfully perturbs more images to cause misclassifications compared to SimBA-DCT (97.8\% success rate). However, SimBA-DCT's slightly better test accuracy may reflect its lower success rate and occasional failures to find adversarial perturbations in the restricted low-frequency DCT basis, as noted in the paper. The trade-off between query efficiency (favoring SimBA-DCT) and attack success (favoring SimBA in pixel space) is evident, and these results highlight the need for further investigation into optimizing the orthonormal basis selection to balance success rate and model robustness.

% Clearing the page
\clearpage
\section{Defenses}


\subsection{Adverserial Training}
\paragraph{Introduction}
Adversarial training is a principled and widely adopted defense mechanism in the field of adversarial machine learning. It aims to improve the robustness of deep neural networks by explicitly incorporating adversarial examples—inputs intentionally perturbed to mislead the model—into the training process. The core idea is to expose the model to these adversarial inputs during training so that it learns to correctly classify both clean and perturbed data, thereby developing resistance to a wide range of potential attacks.

From a mathematical perspective, adversarial training formulates the learning problem as a min-max optimization task. The inner maximization identifies the most damaging perturbation within a constrained set (usually defined by an ℓₚ-norm ball), while the outer minimization updates the model parameters to minimize the loss under these worst-case perturbations. This process effectively strengthens the decision boundaries of the model, making it less sensitive to small but malicious input variations.

A key strength of adversarial training lies in its ability to instill intrinsic robustness into the model, as opposed to post hoc defenses that may be bypassed with adaptive attacks. Despite its computational intensity, adversarial training is considered one of the most effective defenses, especially against white-box attacks where the adversary has full knowledge of the model. Numerous empirical studies have demonstrated that models trained with adversarial examples outperform standard models in terms of adversarial robustness and generalization to unseen attack strategies.
\cite{zheng2020efficientadversarialtrainingtransferable}

\paragraph{Methodology}

\subparagraph{Ways of Adversarial Training} 

 Adversarial training strategies differ significantly based on how adversarial examples are generated and incorporated during training. Two primary approaches are \textbf{fixed adversarial training} and \textbf{on-the-fly adversarial training}, each with distinct implications for robustness and efficiency.
\begin{itemize}
    \item \textbf{Fixed Adversarial Training}
    \begin{itemize}
        \item Adversarial examples are pre-computed before the training process begins and remain unchanged throughout training.
        \item The same set of perturbed samples is reused across all training epochs.
        \item As the model evolves, these fixed examples may no longer pose a challenge, reducing their effectiveness.
        \item The model may overfit to specific perturbations in the fixed set, learning to resist only those patterns without achieving general robustness.
        \item This often results in poor generalization to novel or adaptive adversarial attacks.
    \end{itemize}

    \item \textbf{On-the-Fly Adversarial Training}
    \begin{itemize}
        \item Adversarial examples are generated dynamically during each training iteration using the current state of the model.
        \item The perturbations evolve alongside the model, ensuring that they remain challenging and relevant throughout training.
        \item This method helps the model develop more general and adaptive robustness against a wide range of attacks.
        \item It discourages memorization of specific attack patterns and promotes stronger decision boundaries.
        \item Although computationally expensive (up to 7× to 100× more than standard training), techniques such as \textbf{DEAT}(Dynamic Efficient Adversarial Training) and \textbf{ATTA} (Adversarial Training with Transferable Adversarial Examples) reduce training costs while preserving robustness.
    \end{itemize}
\end{itemize}

\subparagraph*{Fine-Tuning Phase}

The defensive mechanism employed during the training phase is based on an adaptive adversarial training strategy. In each training step, the input batch is divided into two parts: a \textbf{clean portion} and an \textbf{adversarial portion}. The ratio between these portions gradually transitions from 75\% clean to 50\% clean as training progresses, allowing the model to be increasingly exposed to adversarial perturbations over time.

Adversarial examples are generated \textbf{on-the-fly} using either the FGSM or PGD attack, selected probabilistically based on the training epoch. Early in training, the FGSM method is preferred due to its efficiency, while in later epochs, PGD becomes more likely to ensure stronger adversarial challenges. For each adversarial example, the perturbation magnitude $\epsilon$ is randomly selected from a predefined and epoch-scaled list. The PGD attack further adjusts the number of iterations dynamically based on the training progress.

Once clean and adversarial examples are prepared, they are concatenated into a single batch. The model is then trained on this combined dataset in a single forward-backward pass, updating its parameters accordingly. This dynamic and progressively challenging training setup encourages the model to learn robust decision boundaries, providing effective defense against both weak and strong adversarial attacks.
\subparagraph*{Evaluation Procedure During Test Phase}
To evaluate the effectiveness of adversarial training, we assess model performance on adversarial examples generated from the test set, which contains 512 samples. Prior to applying adversarial training, test-time adversarial examples are crafted using the FGSM and PGD attack methods and evaluated against the baseline (undefended) model to measure initial vulnerability. Following the adversarial fine-tuning phase, the same attack algorithms are used again to generate new adversarial examples, this time targeting the fine-tuned (defended) model. This ensures a fair and consistent comparison. The accuracy on these post-defense adversarial examples is then reported, highlighting the improvement achieved through adversarial training as a defense mechanism.


\clearpage
\paragraph{Results}\mbox{}\\[1ex]

\begin{table}[H]
\centering
\caption{FGSM Accuracy Before and After Adversarial Training}
\begin{tabular}{ccc}
\toprule
\textbf{Epsilon ($\varepsilon$)} & \textbf{Before Adversarial Training (\%)} & \textbf{After Adversarial Training (\%)} \\
\midrule
0.007 & 44.96 & 77.91 \\
0.010 & 34.30 & 75.39 \\
0.030 & 12.79 & 75.78 \\
0.100 & 9.30  & 47.67 \\
\bottomrule
\end{tabular}
\label{tab:fgsm-training-results}
\end{table}
\begin{table}[H]
\caption{PGD Accuracy Before and After Adversarial Training}
\hspace*{-1cm}% Adds a negative shift to move left
\makebox[0.8\textwidth][l]{ % Reduced width to 0.8\textwidth
\begin{tabular}{cccc}
\toprule
\textbf{Epsilon ($\varepsilon$)} & \textbf{Alpha ($\alpha$)} & \textbf{Before Adversarial Training (\%)} & \textbf{After Adversarial Training (\%)} \\
\midrule
0.007  & 0.00175 & 36.43 & 68.60 \\
0.010  & 0.00250 & 21.71 & 61.43 \\
0.030  & 0.00750 & 2.71  & 42.83 \\
0.100  & 0.02500 & 0.00  & 13.76 \\
\bottomrule
\end{tabular}
}
\label{tab:pgd-training-results}
\end{table}
\paragraph{Discussion of Results}\mbox{}\\[1ex]
The results in Tables~\ref{tab:fgsm-training-results} and~\ref{tab:pgd-training-results} clearly demonstrate the efficacy of the adaptive adversarial training strategy. Across all tested perturbation levels, substantial improvements in model robustness are observed post-training. For FGSM attacks, the defended model maintains accuracy above 75\% for $\varepsilon \leq 0.03$, compared to severe degradation in the undefended baseline (as low as 12.79\%). Even under stronger perturbations ($\varepsilon = 0.1$), the adversarial training yields a notable 38.37 percentage point improvement.

Similarly, under PGD attacks---which are more challenging due to their iterative nature---the trained model shows significant resilience. Notably, at $\varepsilon = 0.03$, accuracy improves from  2.71\% to 42.83\%, a nearly \textbf{16-fold relative increase}. While performance under the strongest PGD setting ($\varepsilon = 0.1$) remains modest (13.76\%), the improvement from 0\% baseline confirms that the defense still imparts non-trivial robustness in extreme cases.

These gains validate the effectiveness of gradually intensifying exposure to adversarial perturbations during training. By shifting from FGSM to PGD and adjusting the clean-to-adversarial ratio over time, the model successfully learns more robust feature representations capable of resisting both single-step and iterative attacks.

\clearpage

\subsection{Noise Fusion}
\paragraph{Introduction}
The Noise-Fusion Method (NFM) is a straightforward and robust defense mechanism crafted to safeguard deep neural network (DNN) image classification models against adversarial attacks. Adversarial attacks exploit DNN vulnerabilities by introducing subtle, often imperceptible, perturbations to input data, such as images, to induce misclassifications. Unlike many sophisticated defense strategies that demand detailed knowledge of the model architecture or specifics of the attack, NFM operates agnostically, requiring no such information, which enhances its versatility and applicability. The core of NFM lies in strategically adding controlled noise—specifically Uniform, Gaussian, or Poisson noise—to both the input data during runtime and the training data during model training. This dual noise injection serves a twofold purpose: it masks adversarial perturbations by "drowning" them within the added noise, rendering the attack less effective, and it bolsters the model’s resilience to noisy inputs, thereby improving overall robustness. NFM has been rigorously evaluated against prominent adversarial attacks, including the Fast Gradient Signed Method (FGSM), Projected Gradient Descent (PGD), and Sparse L1 Descent (SLD), demonstrating significant improvements in model accuracy under attack conditions. By perturbing attack gradients and stabilizing model performance through noise, NFM offers a generalized defense with potential applicability beyond image-based models to other data types, such as audio, making it a promising approach for securing DNNs against a wide range of adversarial threats without the need for attack-specific or model-specific tailoring.
\cite{shi2022noisefusion}
\paragraph{Methodology}



\subparagraph{Noise Fusion Mechanism}

Let $\mathbf{x} \in \mathbb{R}^{H \times W \times C}$ denote an input image and $\mathbf{n} \in \mathbb{R}^{H \times W \times C}$ a noise matrix sampled from a specific distribution. The fused input $\mathbf{x}_{\text{noise}}$ is computed as:

\begin{equation}
    \mathbf{x}_{\text{noise}} = (1 - \alpha) \cdot \mathbf{x} + \alpha \cdot \mathbf{n}, \quad \alpha \in [0, 1]
\end{equation}

Here, $\alpha$ is the noise amplitude, which controls the strength of the noise blended into the image.
\subparagraph*{Noise Distributions}

The following three noise distributions are used to evaluate and enhance model robustness:

\begin{itemize}
    \item \textbf{Uniform Noise}: Each pixel is sampled independently from a uniform distribution over the full intensity range:
    \begin{equation}
        \mathbf{n}_{i,j,c} \sim \mathcal{U}(0, 255)
    \end{equation}
    
    \item \textbf{Gaussian Noise}: Each pixel is sampled from a Gaussian distribution centered at the midpoint of the valid pixel range:
    \begin{equation}
        \mathbf{n}_{i,j,c} \sim \mathcal{N}(\mu, \sigma^2), \quad \mu = 127.5,\ \sigma = 64
    \end{equation}
    
    \item \textbf{Poisson Noise}: Poisson noise is signal-dependent and is drawn per pixel based on the input intensity:
    \begin{equation}
        \mathbf{n}_{i,j,c} \sim \text{Poisson}(\lambda = \mathbf{x}_{i,j,c})
    \end{equation}
\end{itemize}

\subparagraph{Training and Inference Procedure}

Noise fusion is applied in both training and evaluation as follows:

\textbf{Training Phase:}
\begin{enumerate}
    \item Sample noise $\mathbf{n} \sim \mathcal{D} \in \{\text{Uniform}, \text{Gaussian}, \text{Poisson}\}$
    \item Compute $\mathbf{x}_{\text{noise}} = (1 - \alpha) \cdot \mathbf{x} + \alpha \cdot \mathbf{n}$
    \item Train the model using $(\mathbf{x}_{\text{noise}}, y)$
\end{enumerate}

\textbf{Inference Phase:}
\begin{enumerate}
    \item Sample noise $\mathbf{n} \sim \mathcal{D}$
    \item Compute $\mathbf{x}_{\text{noise}} = (1 - \alpha) \cdot \mathbf{x} + \alpha \cdot \mathbf{n}$
    \item Pass $\mathbf{x}_{\text{noise}}$ through the model to obtain predictions
\end{enumerate}

\textbf{Comparison Between Defense Phases}

\begin{table}[H]
\centering
\begin{tabular}{|c|l|l|}
\hline
\textbf{Phase} & \textbf{Operation} & \textbf{Purpose} \\
\hline
Training  & Add noise to images & Improve robustness via augmentation \\
Inference & Fuse noise into input & Suppress adversarial perturbations \\
\hline
\end{tabular}
\caption{Noise Fusion Strategy Across Phases}
\end{table}



% Results section
\paragraph{Results}\mbox{}\\[1ex]



\begin{table}[ht]
\centering
\scriptsize
\caption{Test-set accuracies before and after Noise Fusion, with absolute improvement}
\label{tab:noise-fusion-accuracies}
\renewcommand{\arraystretch}{1.5} % <-- Increases row height
\resizebox{\textwidth}{!}{%
\begin{tabular}{@{}l l c c c c c@{}}
\toprule
\textbf{Attack Type} & \textbf{Noise Distribution} & \boldmath$\epsilon$ & \textbf{Alpha} & \textbf{Test Accuracy after Attack (\%)} & \textbf{Test Accuracy after Noise Fusion (\%)} & \textbf{Improvement (\% )} \\
\midrule
FGSM & Gaussian & 0.007 & --      & 58.00 & 95.16 & 37.16 \\
FGSM & Gaussian & 0.01  & --      & 44.00 & 93.41 & 49.41 \\
FGSM & Gaussian & 0.03  & --      & 14.00 & 83.33 & 69.33 \\
FGSM & Gaussian & 0.10  & --      &  6.00 & 48.64 & 42.64 \\
PGD  & Gaussian & 0.007 & 0.00175 & 35.00 & 93.22 & 58.22 \\
PGD  & Gaussian & 0.01  & 0.00250 &  7.00 & 80.23 & 73.23 \\
PGD  & Gaussian & 0.03  & 0.00750 &  3.00 & 67.05 & 64.05 \\
PGD  & Gaussian & 0.10  & 0.02500 &  0.00 & 42.83 & 42.83 \\
\midrule
FGSM & Uniform  & 0.007 & --      & 58.00 & 83.72 & 25.72 \\
FGSM & Uniform  & 0.01  & --      & 44.00 & 81.59 & 37.59 \\
FGSM & Uniform  & 0.03  & --      & 14.00 & 72.48 & 58.48 \\
FGSM & Uniform  & 0.10  & --      &  6.00 & 42.64 & 36.64 \\
PGD  & Uniform  & 0.007 & 0.00175 & 35.00 & 68.99 & 33.99 \\
PGD  & Uniform  & 0.01  & 0.00250 &  7.00 & 39.92 & 32.92 \\
PGD  & Uniform  & 0.03  & 0.00750 &  3.00 & 25.39 & 22.39 \\
PGD  & Uniform  & 0.10  & 0.02500 &  0.00 & 10.66 & 10.66 \\
\midrule
FGSM & Poisson  & 0.007 & --      & 58.00 & 93.99 & 35.99 \\
FGSM & Poisson  & 0.01  & --      & 44.00 & 92.44 & 48.44 \\
FGSM & Poisson  & 0.03  & --      & 14.00 & 76.55 & 62.55 \\
FGSM & Poisson  & 0.10  & --      &  6.00 & 42.83 & 36.83 \\
PGD  & Poisson  & 0.007 & 0.00175 & 35.00 & 46.12 & 11.12 \\
PGD  & Poisson  & 0.01  & 0.00250 &  7.00 &  7.36 &  0.36 \\
PGD  & Poisson  & 0.03  & 0.00750 &  3.00 &  2.91 & -0.09 \\
PGD  & Poisson  & 0.10  & 0.02500 &  0.00 &  0.97 &  0.97 \\
\bottomrule
\end{tabular}%
}
\end{table}
\clearpage
As shown in Table~\ref{tab:noise-fusion-accuracies}, the Noise Fusion Method significantly improves model robustness against adversarial attacks, with \textbf{Gaussian noise} demonstrating the strongest and most consistent defense across both FGSM and PGD attacks. Specifically, Gaussian noise restores accuracy by up to \textbf{+73.23\%} under PGD and \textbf{+69.33\%} under FGSM, outperforming Uniform and Poisson noise distributions. While Uniform noise offers moderate gains, especially under FGSM, its effectiveness diminishes with stronger attacks like PGD. Poisson noise shows good performance against FGSM but fails to defend against PGD, offering minimal or even negative improvements at higher $\epsilon$ values. Overall, Gaussian noise emerges as the most robust and reliable choice for noise fusion, particularly as attack intensity increases.


\subsection{JPEG Compression Defense}
\label{sec:jpeg_defense}

\paragraph{Introduction}
JPEG compression is a widely used lossy image compression standard that can also serve as a simple yet effective preprocessing defense against adversarial attacks \cite{das2017keepingbadguysout}. The core idea is that adversarial perturbations, especially those generated by gradient-based methods like FGSM and PGD, are often composed of high-frequency noise finely tuned to the model's sensitivities. The quantization step in the JPEG compression algorithm is particularly effective at discarding these high-frequency components. By compressing and then decompressing an adversarial image, the carefully crafted perturbation can be disrupted or "smoothed" out, potentially restoring the image's original classification. This makes JPEG compression a practical, model-agnostic defense strategy that can be applied as a preprocessing step before inference.

\paragraph{Results and Analysis}
To quantify the efficacy of this defense, we evaluated its impact on adversarial examples generated by both FGSM and PGD attacks across a range of epsilons. We applied four standard JPEG quality levels: 25 (high compression, low quality), 50, 75, and 95 (low compression, high quality). The comprehensive results are presented in Table~\ref{tab:jpeg_defense_merged}.

The data reveals two clear trends across both attack types. First, there is a strong inverse correlation between JPEG quality and defensive strength. The most aggressive compression setting (Q=25) consistently provided the greatest improvement in accuracy**. For instance, against a PGD attack with $\epsilon=0.01$, JPEG compression at Q=25 restored accuracy from 21.71\% to 68.80\% (a +47.09\% improvement), whereas the mildest compression at Q=95 only restored it to 48.84\%. This indicates that the higher distortion introduced by low-quality compression is more effective at destroying the adversarial noise.

Second, the defense's utility diminishes as the attack strength ($\epsilon$) increases. While the defense provides substantial gains for small epsilons, its effectiveness wanes against the very strong perturbations of $\epsilon=0.1$. Against the PGD attack at this level, accuracy was restored from 0\% to only 8.91\% at Q=25, and a mere 0.78\% at Q=95. This suggests that when the adversarial perturbation is large and coarse, it is robust enough to survive the compression process.

The results also allow for a direct comparison between FGSM and PGD. While PGD is a stronger attack that reduces the initial adversarial accuracy to lower levels, the JPEG defense restores accuracy to a comparable level for both attacks at a given $\epsilon$. For example, at $\epsilon=0.03$ and Q=25, the defended accuracy is 38.18\% for FGSM and 38.37\% for PGD, demonstrating the defense's consistent performance against both single-step and iterative gradient-based attacks.

\begin{table}[H]
    \centering
    \caption{Model Accuracy Against FGSM and PGD Attacks with JPEG Compression Defense.}
    \label{tab:jpeg_defense_merged}
    \begin{tabular}{cc|cc|cc|cc}
    \toprule
    \multirow{2}{*}{\textbf{$\epsilon$}} & \multirow{2}{*}{\textbf{Quality}} & \multicolumn{2}{c|}{\textbf{Adversarial Acc.}} & \multicolumn{2}{c|}{\textbf{Defended Acc.}} & \multicolumn{2}{c}{\textbf{Improvement}} \\
    \cmidrule{3-8}
    & & \textbf{FGSM} & \textbf{PGD} & \textbf{FGSM} & \textbf{PGD} & \textbf{FGSM} & \textbf{PGD} \\
    \midrule
    \multirow{4}{*}{0.007} & 25 & 44.96 & 36.43 & 81.40 & 79.26 & +36.43 & +42.83 \\
    & 50 & 44.96 & 36.43 & 80.04 & 75.58 & +35.08 & +39.15 \\
    & 75 & 44.96 & 36.43 & 78.88 & 73.26 & +33.91 & +36.82 \\
    & 95 & 44.96 & 36.43 & 71.90 & 63.18 & +26.94 & +26.74 \\
    \midrule
    \multirow{4}{*}{0.010} & 25 & 34.30 & 21.71 & 72.87 & 68.80 & +38.57 & +47.09 \\
    & 50 & 34.30 & 21.71 & 69.96 & 64.53 & +35.66 & +42.83 \\
    & 75 & 34.30 & 21.71 & 68.80 & 62.21 & +34.50 & +40.50 \\
    & 95 & 34.30 & 21.71 & 56.40 & 48.84 & +22.09 & +27.13 \\
    \midrule
    \multirow{4}{*}{0.030} & 25 & 12.79 & 2.71 & 38.18 & 38.37 & +25.39 & +35.66 \\
    & 50 & 12.79 & 2.71 & 35.08 & 31.98 & +22.29 & +29.26 \\
    & 75 & 12.79 & 2.71 & 29.84 & 25.39 & +17.05 & +22.67 \\
    & 95 & 12.79 & 2.71 & 21.12 & 10.47 & +8.33 & +7.75 \\
    \midrule
    \multirow{4}{*}{0.100} & 25 & 9.30 & 0.00 & 12.02 & 8.91 & +2.71 & +8.91 \\
    & 50 & 9.30 & 0.00 & 10.66 & 6.40 & +1.36 & +6.40 \\
    & 75 & 9.30 & 0.00 & 10.27 & 3.49 & +0.97 & +3.49 \\
    & 95 & 9.30 & 0.00 & 10.08 & 0.78 & +0.78 & +0.78 \\
    \bottomrule
    \end{tabular}
    
\end{table}



\end{document}